\documentclass[a4paper,fleqn]{cas-sc}

\usepackage[authoryear,round]{natbib}
\usepackage{subcaption}
\usepackage{comment}
\usepackage{amsmath}
\usepackage{mathtools}
\DeclareMathOperator*{\argmax}{\arg\,max}
\usepackage[normalem]{ulem} % Added by SamB for the s\out command

\newtheorem{theorem}{Theorem}
\newtheorem{lemma}[theorem]{Lemma}
\newtheorem{corollary}[theorem]{Corollary}
\newtheorem{proposition}{Proposition}
\newtheorem{observation}{Observation}
\newdefinition{definition}{Definition}
\newdefinition{conjecture}{Conjecture}
\newproof{proof}{Proof}

%\numberwithin{equation}{section}

\def\tsc#1{\csdef{#1}{\textsc{\lowercase{#1}}\xspace}}
\tsc{SB}
\tsc{SD}
\tsc{GM}
\tsc{DP}

\newcommand{\samd}[1]{\textcolor{red}{\textbf{#1}}}
\newcommand{\samb}[1]{\textcolor{blue}{\textbf{#1}}}
\newcommand{\dave}[1]{\textcolor{violet}{\textbf{#1}}}
\newcommand{\greta}[1]{\textcolor{teal}{\textbf{#1}}}

\begin{document}

\let\WriteBookmarks\relax
\def\floatpagepagefraction{1}
\def\textpagefraction{.001}
\shorttitle{}
\shortauthors{S. Doak et~al.}

%\title[mode = title]{Dynamics, bifurcations and asymptotics in unified growth models}
\title[mode = title]{Dynamics, bifurcations and predictions of unified growth models}

\author[1]{Sam Doak}[type=editor]
\credit{Conceptualization, Methodology, Formal analysis, Writing - Original draft, Writing - Review \& Editing}

\author[2]{Greta Meggiorini}[type=editor, orcid=0000-0001-5098-9740]
\credit{Conceptualization, Methodology, Validation, Writing - Original draft, Writing - Review \& Editing}

\cormark[1]
\cortext[cor1]{Corresponding author}
\ead{greta.meggiorini@auckland.ac.nz}

\author[1]{Samuel Bolduc-St-Aubin}[type=editor]
\credit{Investigation, Formal analysis, Visualization, Writing - Original draft, Writing - Review \& Editing}

\author[1]{Davide Papapicco}[type=editor, orcid=0000-0002-5091-2599]
\credit{Project administration, Conceptualization, Formal analysis, Writing - Original draft, Writing - Review \& Editing}

\affiliation[1]{organization={Department of Mathematics, University of Auckland},
                 city={Auckland},
                 country={New Zealand}}
\affiliation[2]{organization={Department of Economics, University of Auckland},
                 city={Auckland},
                 country={New Zealand}}
\begin{titlepage}

\begin{abstract}
Lorem ipsum dolor sit amet, consectetur adipiscing elit, sed do eiusmod tempor incididunt ut labore et dolore magna aliqua. Ut enim ad minim veniam, quis nostrud exercitation ullamco laboris nisi ut aliquip ex ea commodo consequat. Duis aute irure dolor in reprehenderit in voluptate velit esse cillum dolore eu fugiat nulla pariatur. Excepteur sint occaecat cupidatat non proident, sunt in culpa qui officia deserunt mollit anim id est laborum.
\end{abstract}

% % LEAVE THIS HERE FOR SUBMISSION
%\begin{graphicalabstract}
% %\includegraphics{figs/cas-grabs.pdf}
% %\end{graphicalabstract}
% %\begin{highlights}
% %\item Research highlights item 1
% %\item Research highlights item 2
% %\item Research highlights item 3
% %\end{highlights}

\begin{keywords}
Long-run growth \sep 
Dynamical systems \sep
Malthusian trap \sep
Bifurcation analysis \sep
\end{keywords}

\end{titlepage}

\maketitle

\section{Introduction}\label{sec:introduction}

Unified growth theory studies the transition from Malthusian stagnation to sustained modern-growth within a single dynamic framework. The objective is to explain why economies can spend long periods with low income per capita, slow technological progress, high fertility, and little investment in human capital, before entering a regime characterized by sustained technological growth, rising education, and declining fertility. A central idea is that population, technology, human capital, and fertility are jointly determined over the course of development \citep{Galor2000,Galor2005,Galor2011}.

The unified-growth literature builds on earlier work linking population and technological change over the very long run \citep{Kremer1993} and has developed alternative mechanisms for the transition from stagnation to growth, including human-capital accumulation, fertility decline, life expectancy, and sectoral transformation \citep{GalorMoav2002,HansenPrescott2002,Doepke2004,CervellatiSunde2005,StrulikWeisdorf2008}.

A canonical contribution is \citet{Galor2000}, who develop a mechanism in which population affects technological progress, technological progress raises the return to human-capital investment, and households adjust fertility and education choices as the economy develops. This mechanism generates the three stages emphasized by unified growth theory: the Malthusian stagnation (often referred to as the Malthusian trap), the post-Malthusian regime, and the modern-growth regime (often referred to as the sustained growth state). 

The Galor--Weil model (henceforth, GW) uses discrete time dynamical systems to describe these three states and the transitions connecting them. 
In particular, the GW model specifies, under the unified growth theory framework, generic conditions on endogenous state variables to constitute a qualitative family of discrete time dynamical systems.
To assess the capability of the GW model to accurately capture the transitions of economies out of the Malthusian stagnation and onto sustained growth, one needs to specify functional forms that extract a specific dynamical system out of the qualitative family.
This is done in \citet{Lagerlof2006} who shows that, under a particular set of parameters and a particular starting state, an economy can transition out of the Malthusian trap and onto a state of sustained growth.

This paper builds upon these works and asks whether this transition path is a generic prediction of the quantitative Galor--Weil--Lagerl{\"o}f model (henceforth, GWL) or an observation of the specified parameter values and starting conditions.
While a single simulated trajectory can illustrate an economic mechanism, it does not suffice to characterize the behavior of the complex evolution of global economies. In models with subsistence constraints, threshold effects, and piecewise-defined policy functions, different initial conditions or parameter values may generate qualitatively different outcomes.

The overarching aim of this paper is thus to show how, using well-established tools from dynamical systems and bifurcation theory, one can uncover all the qualitatively different trajectories of an economy starting from any initial condition and for any choice of its parameter values.
By applying these standard mathematical tools to the GWL model, our hope is to introduce a framework for the modeling, analysis and forecasting of complex, long-run macroeconomic behaviors across different societies.
To showcase the efficacy of this analytical framework, we obtain a series of results concerning the dynamics of the GWL model, which we list below.

First, we demonstrate that the family of four-dimensional dynamical systems of the GW model is only three-dimensional. In \citet{Galor2000}, the state of the economy is characterized by technological growth, education, technology, and population. However we show that, along equilibrium paths, education is not an independent state variable but rather a policy function of technological growth. As such, the relevant endogenous state vector contains technological growth, technology, and population, while education, human capital, potential income, and fertility are determined from the state. This is pivotal for the development of unified growth theory, as a three-dimensional system is significantly less complex to study than a four-dimensional one.
 
Second, the transition path from Malthusian stagnation to sustained growth is not a global property of the GWL model. We identify open regions of the state space in which economies experience population collapse within one or two generations. These collapses arise from the interaction between the subsistence constraint, fertility choices, and the dependence of potential income on technology and population. Thus, economies with initial conditions nearby to that specified in \citet{Lagerlof2006} can generate sharply different long-run outcomes within the same calibrated model.
    
Third, the only steady-states available to an economy in the GWL model are those of population collapse. Specifically, neither the Malthusian trap nor the sustained growth state, are valid steady-states of the GWL model. With respect to the sustained growth state, observed in the benchmark simulation of \citet{Lagerlof2006}, we show that it can become a steady-state of the GWL system only in a two-dimensional asymptotic reduction where the level of technology is high. In this regime, the potential income is large enough so that the subsistence constraint no longer binds, the dynamics of technological growth and population decouple from the that of the level of technology, thus reducing the system from three to two dimensions.

Fourth, the dynamics in the asymptotic regime is organized by switching manifolds, generated by the education threshold, the population scale-effect threshold, and the fertility rule. These manifolds partition the state space into regions with distinct local dynamics. We show that, under Lagerl\"{o}f's calibration of the parameters, there exist infinitely many fixed points in the modern-growth regime. Using bifurcation analysis, we also show that such result is not robust. If one perturbs the preference weight, which governs the household's valuation of children relative to adult consumption, such steady-states either collapse to a single one or they disappear altogether.

This paper contributes to the mathematical development of unified growth theory by turning a calibrated historical simulation into a global characterization of model predictions.
In particular, it connects to the broader literature on nonlinear dynamics in economic models, where global behavior, invariant sets, and cycles are central for the understanding of different equilibrium outcomes \citep{BenhabibDay1982,Matsuyama1999,GomisPorquerasHaro2003}.
This complements a recent quantitative work, in unified growth theory, that studies the ability of nonlinear, long-run development models to match empirical growth patterns \citep{CervellatiMeyerheimSunde2023}.

The rest of the manuscript is organized as follows. Section~\ref{sec:background} reviews the Galor--Weil environment and Lagerl{\"o}f's quantitative implementation.We also provide a brief primer on the mathematical techniques and fundamental notions of dynamical systems used throughout the analysis.  In Section~\ref{sec:analysis} we study the global dynamics of economies evolving according to the GWL model and derive the results listed above. In Section~\ref{sec:smoothing} we revisit the GWL model by prescribing smoother, but still economically relevant, evolution laws for the endogenous and choice variables. Section~\ref{sec:conclusions} discusses our findings and details future directions for the analysis  and predictions of unified growth models.

\section{Background}\label{sec:background}
This section reviews the unified growth framework of \citet{Galor2000} and the quantitative implementation of \citet{Lagerlof2006}. In the following exposition we separate the production side of the economy from the household optimization problem and obtain the resulting laws of motion.

\subsection{The GW model}\label{subsec:galor_weil_environment}
The GW model describes deterministic, overlapping-generations economies evolving in discrete time. Members of generation $t$ live for two periods: they are children in period $t-1$ and adults in period $t$. As adults, they earn income, consume, choose fertility, and invest in the education of children who become the adults of generation $t+1$.

\subsubsection{Economic production and potential income}\label{subsubsec:production}
Following \citet{Galor2000} and \citet{Galor2011}, under the simplifying assumptions of a fixed and normalized supply of land, and no property rights over such land, an adult agent in period $t$ can earn a maximum potential income 
%
\begin{equation}\label{eq:income_equality} 
    z_t = h_t^{\alpha} \bigg(\frac{A_t}{L_t}\bigg)^{1-\alpha},
\end{equation} 
%
where $A_t$ is the level of technology in the economy, $L_t$ is the population of adult agents, $h_t$ is the human capital of the agents, and $\alpha\in(0,1)$ is the labor share. 
The human capital of the adult generation $t$ (which are children in generation $t-1$) is given by 
%
\begin{equation}\label{eq:human_capital} 
    h_{t} \coloneqq h(e_{t},G_{t}), 
\end{equation} 
%
where $e_{t}\geq 0$ is the adults' choice for the education of children in generation $t-1$, and 
%
\begin{equation}\label{eq:relative_increase_technology} 
    G_{t} \coloneqq \frac{A_{t}-A_{t-1}}{A_{t-1}} 
\end{equation} 
%
is the rate of technological growth between generations $t-1$ and $t$.

\subsubsection{Optimal preferences and the household problem}\label{subsubsec:household_problem}
Adult agents of generation $t$ determine their preferred number of children $n_t\geq 0$ and the education $e_{t+1}\geq 0$ to provide to them, given the state of the economy.
Specifically, at each generation $t$, agents allocate their time between economic production and child-rearing activities according to their preferences, which depend on the endogenous states variables.
In the GW model, these preferences are encoded as a utility function
%
\begin{equation}\label{eq:utility_function} 
    u(n_t,e_{t+1}) = c(n_t,e_{t+1})^{1-\gamma} \big(n_t h(e_{t+1},G_{t+1})\big)^{\gamma}, 
\end{equation} 
%
which represents the trade-off, parametrized by $\gamma \in (0, 1)$, between consumption and the human capital of their children as in \cite{Galor2011}.
% \footnote{This formulation is equivalent to \citet{Galor2005}.} % for the household problem, since the future wage per efficiency unit is taken as given by households and thus drops out of the fertility and education choices. 
The consumption of the agents is provided as 
\begin{equation}\label{eq:consumption_normalized} 
    c(n_t,e_{t+1}) = z_t \big(1-n_t(\tau+e_{t+1})\big). 
\end{equation}
% \begin{equation}\label{eq:consumption} 
%     c(n_t,e_{t+1}) = z_t \big(1-n_t(\tau_q+\tau_e e_{t+1})\big),
% \end{equation} 
%
where $\tau + e_{t+1}$ is the time cost of raising one child, with $\tau$ being a fixed time cost and the variable time cost of education is set to one. 
The household problem is thus obtained by the maximization of the utility of agents of generation $t$ in terms of their choices of fertility $n_t$ and bestowed education $e_{t+1}$. 
Mathematically, the optimal choices $n_t^*$ and $e^*_{t+1}$ solve the following constrained optimization problem
%
\begin{align}\label{eq:optimization_problem}
     &\text{find}\;\;(n_t^*,\,e_{t+1}^*) = \argmax\!\big(\log u(n_t,\,e_{t+1})\big)\;\;\text{subject to}\nonumber\\
     &n_t^*,\,e_{t+1}^*\geq0\;\;\text{and}\;\;c(n_t^*,\,e_{t+1}^*)\geq\tilde{c}\,,
\end{align}
%
where $\tilde{c}>0$ is a cost-of-living quantity identifying the subsistence consumption below which agents cannot survive.

\begin{proposition}\label{prop:optimal_solution}
    Suppose that $z_t\geq \tilde c$, so that the household problem is feasible. The optimal education is given by a function of technological progress
    %
    \begin{equation}\label{eq:optimal_et}
        e_{t+1}^*
        \eqqcolon
        e(G_{t+1})
        =
        \max\!\big(0,\varepsilon(G_{t+1})\big),
    \end{equation}
    %
    where the positive-education branch $\varepsilon_{t+1} = \varepsilon(G_{t+1})$ satisfies
    %
    \begin{equation}\label{eq:implicit_optimal_education}
        \Phi(\varepsilon_{t+1},G_{t+1})
        \coloneqq
        (\tau+\varepsilon_{t+1})
        \frac{\partial h(\varepsilon_{t+1},G_{t+1})}{\partial \varepsilon_{t+1}}
        -
        h(\varepsilon_{t+1},G_{t+1})
        =
        0.
    \end{equation}
    The optimal fertility satisfies
    %
    \begin{equation}\label{eq:optimal_nt}
        n_t^*
        =
        \begin{cases}
            \dfrac{\gamma}{\tau+e_{t+1}^*},
            &
            \dfrac{\tilde c}{1-\gamma}\leq z_t,
            \\[10pt]
            \dfrac{1-\frac{\tilde c}{z_t}}{\tau+e_{t+1}^*},
            &
            \tilde c\leq z_t<\dfrac{\tilde c}{1-\gamma},
        \end{cases}.
    \end{equation}
    %
\end{proposition}

\begin{proof}
    See Appendix \ref{appenix_A}. 
    % The first-order condition for education is obtained by maximizing $h(e_{t+1},G_{t+1})/(\tau+e_{t+1})$. Under the maintained assumptions on $h$, the positive-education branch is locally characterized by the implicit function theorem.
\end{proof}
%
Proposition~\ref{prop:optimal_solution} characterizes the household optimum whenever subsistence consumption is attainable. If the potential income $z_t$ is sufficiently high, the subsistence constraint does not bind and agents devote a constant share $\gamma\in(0,1)$ of time to child-rearing activities. 
If instead the potential income $z_t$ is above subsistence and below $\tilde c/(1-\gamma)$, the subsistence constraint binds and fertility rises with potential income. 
Finally, if the potential income is below subsistence, no feasible allocation can deliver subsistence consumption. Following the demographic closure used in \citet{Lagerlof2006}, fertility is set to $n_t^*=0$. Whenever this happens, the economy undergoes population collapse, i.e. $L_{t+1}=0$, meaning that the population of agents in generation $t+1$ goes extinct.

\subsubsection{Derivation of the laws of motion}\label{subsubsec:law_of_motion} 
\samb{An important comment which we must discuss: the notation for $e_t$ and $e(G_t)$ is interchanged constantly throughout the text. We must choose ONE notation and stick to it consistently. I believe we should keep the $e_t$ notation from Eq.(12) onward and stop using $e(G_t)$.}

In order to derive the laws of motion for an economy described by unified-growth theory we need to determine how the relevant variables evolve over time. Note that the solution to the household problem prescribes precisely this for $e_t$ and $L_t$. Moreover, technological progress $G_{t+1}$, defined in \eqref{eq:relative_increase_technology}, is specified as a function $G_{t+1} = \mathcal G(e_t,L_t)$ of education $e_t$ of children of generation $t$ and the adult agents population $L_t$.
We remark that, within the GW model, the specific form of $\mathcal G$ is left unspecified, subject to the conditions in Table~\ref{table:conditions}. 
Combining the optimum of the household problem \eqref{eq:optimal_nt}--\eqref{eq:optimal_et}, the definition of technological growth \eqref{eq:relative_increase_technology}, \citet{Galor2000} derive a family of dynamical systems in discrete time given as 
%
\begin{align}\label{eq:galor_formulation} 
    \begin{cases} 
        G_{t+1} &= \mathcal G(e_t,L_t),\\
        e_{t+1} &= e(G_{t+1}) = \max\!\big(0,\varepsilon(\mathcal G(e_t,L_t))\big), \\
        A_{t+1} &= (1+G_{t+1})A_t = (1+\mathcal G(e_t,L_t))A_t, \\
        L_{t+1} &= n_t^*L_t. 
    \end{cases} 
\end{align} 
%
In \citet{Galor2000}, and in the literature that followed \citep{Galor2005,Galor2011,Lagerlof2006}, the dynamical system \eqref{eq:galor_formulation} is represented in terms of four variables $(G_t,e_t,A_t,L_t)$. 
However, once the household problem has been solved, education is not an independent state variable.
In particular the education at time $t+1$ depends only on the rate of technological progress at the same time $t+1$.
Hence, by shifting the time index, $e_t$ can be obtained from $G_t$, implying that the current education is an endogenous policy function of the current technological growth and, consequently, does not require an independent law of motion.

The same logic applies to human capital $h_t$, potential income $z_t$, and optimal fertility $n^*_t$.
More specifically, given the state variables $(G_t,A_t,L_t)$ at period $t$,
%
\begin{equation}\label{eq:choice_variables}
    z_t=h(e(G_t),G_t)^{\alpha}
    \bigg(\frac{A_t}{L_t}\bigg)^{1-\alpha},\qquad h_t=h(e(G_t),G_t),\qquad \text{ and }\qquad n_t^*=n(G_t,A_t,L_t)
\end{equation}
are fully determined within the same period by \eqref{eq:income_equality}, \eqref{eq:human_capital} and \eqref{eq:optimal_nt}.
%
% To facilitate the distinction between endogenous state and choice variables, we denote them with upper and lower case, respectively; see, for example, $G_t$ and $e_t$.
% \sout{To facilitate the important distinction between endogenous state variables (e.g. technological progress) and endogenous (pre-determined) choice variables (e.g. education) we will henceforth denote with an uppercase letter the former (e.g. $G_t$) and with a lowercase letter the latter (e.g. $e_t$).}
The dynamics can therefore be represented using the reduced endogenous state vector $X_t=(G_t,A_t,L_t)\in\mathbb R_{\geq0}^3$ with the corresponding laws of motion
%
\begin{align}\label{eq:galor_family}
    \begin{cases}
        G_{t+1} &= \mathcal G(e(G_t),L_t),\\
        A_{t+1} &= \big(1 + G_{t+1}\big)A_t,\\
        L_{t+1} &= n(G_t,A_t,L_t)L_t,
    \end{cases}
\end{align}
%
where $e(G_t)$ is given by \eqref{eq:optimal_et}.

%\subsubsection{Accretion and erosion of human capital and technological growth}\label{subsubsec:family_definition}

%We remark that the reduction from $(G_t,e_t,A_t,L_t)$ to $(G_t,A_t,L_t)$, while removing a redundant variable $e_t=e(G_t)$, does not affect the qualitative equilibrium paths predicted by the GW model. 
While, in fact, \citet{Galor2000} leave the functional forms of the human capital function $h$ (see Eq.~\eqref{eq:human_capital}) and technological growth $\mathcal G$ unspecified, they prescribe a set of conditions they must satisfy. 
In particular, the human capital function $h$ captures a central mechanism in the GW model.
As education raises the human capital of children, technological progress erodes the usefulness of knowledge acquired from the previous generation. This erosion effect is the channel through which technological progress creates an incentive to invest in education.
Similar conditions are specified for the function $\mathcal G$ quantifying the inter-generational rate of technological progress.
We collect these qualitative restrictions for the functional choices of education and technological progress in Table~\ref{table:conditions}.

\begin{table}[width=\linewidth,cols=2,pos=t]
\caption{List of conditions for the unspecified functional forms $\mathcal G(e(G_t), L_t)$ and $h(e(G_t),G_t)$ as provided in \citet{Galor2000}.}\label{table:conditions}
\begin{tabular*}{\tblwidth}{@{} LL@{} }
\toprule
Condition & Interpretation \\
\midrule
$\mathcal{G}(0,\,L_t)\geq0$ & Even in the absence of education, technology can never decrease \\
$\partial_{e}\mathcal{G},\,\partial_{L_t}\mathcal{G}>0$ & Technology progresses faster at higher education and population level \\
$\partial^2_{e}\mathcal{G},\,\partial^2_{L_t}\mathcal{G}<0$ & The effects of accretion of technological progress are marginally diminishing \\
$\partial_{e}h>0,\,\partial_{G_t}h<0$ & Education accrues human capital, while technological progress erodes it\\
$\partial^2_{e}h<0,\,\partial^2_{G_t}h>0$ & The effects of accretion and erosion of human capital are marginally diminishing \\
\bottomrule
\end{tabular*}
\end{table}

\subsection{The GWL model}\label{subsec:lagerlof_exercise}
The framework established by the GW model leaves the human-capital function $h$ and the technological progress function $\mathcal G$ unspecified. 
In order to make predictions about the possible transition paths of economies within unified growth theory, one needs to specify such functions and thereby extract a specific dynamical system out of qualitative family \eqref{eq:galor_family}.
In \citet{Lagerlof2006} such functional forms are assigned and a specific nonlinear map for the endogenous state vector $X_t=(G_t,A_t,L_t)$ is obtained. 
In general, a specific set of laws of motion can be obtained as follows: 
%
\begin{enumerate} 
    \item specify functional forms for $h(e,G)$ and $\mathcal G(e,L)$ that implement the qualitative mechanisms summarized in Table~\ref{table:conditions}; 
    \item use the chosen human capital function $h(e,G)$ to solve the education first-order condition \eqref{eq:implicit_optimal_education} and obtain the education policy function; 
    \item substitute $e_t=e(G_t)$ into the technology equation to obtain the reduced law of motion $G_{t+1} = \mathcal G(e(G_t),L_t)$; 
    \item substitute the education policy function $e(G_t)$ and the human-capital function $h(e(G_t),G_t)$ into the fertility rule \eqref{eq:optimal_nt} to obtain $n(G_t,A_t,L_t)$.
\end{enumerate} 
%
This procedure closes the model while preserving the core mechanism of the GW model: population and education affect technological progress, technological progress affects the return to education, and education and fertility determine the evolution of population and human capital.
Lagerl{\"o}f's choices for the human capital function and the evolution of technological progress are 
%
\begin{equation}\label{eq:lagerlof_choices} 
    h(e_t,G_t) = \frac{e_t+\rho\tau}{e_t+\rho\tau+G_t}, \qquad \mathcal G(e_t,L_t) = (e_t+\rho\tau)\min(\theta L_t,\hat L), 
\end{equation} 
%
which are parametrized by $\rho\in(0,1)$, $\theta>0$, and $\hat L>0$. The term $\min(\theta L_t,\hat L)$ imposes an upper bound on the population scale effect in technological progress. 
Specifically, population raises technological progress up to the threshold 
%
\begin{equation}\label{eq:population_threshold} 
    L = \frac{\hat L}{\theta}, 
\end{equation} 
%
after which the scale effect saturates. 
Substituting the human capital function \eqref{eq:lagerlof_choices} into the first-order condition \eqref{eq:implicit_optimal_education} gives the education policy function 
%
\begin{equation}\label{eq:lagerlof_e} 
    e(G_t) = \max\!\Big(0,-\rho\tau+\sqrt{\tau(1-\rho)G_t}\Big). 
\end{equation} 
%
which is written with time index $t$ after applying the time shift discussed in Section~\ref{subsec:galor_weil_environment}.
Notice how the optimal education vanishes whenever technological progress is sufficiently low; in particular, the optimal education is only positive once $G_t$ crosses the threshold
%
\begin{equation}\label{eq:education_threshold} 
    G = \frac{\rho^2\tau}{1-\rho}. 
\end{equation} 
%
Using \eqref{eq:lagerlof_choices} and \eqref{eq:lagerlof_e}, the GWL model can be written as the map $f:\mathbb R_{\geq0}^3\to\mathbb R_{\geq0}^3$ defined as
%
\begin{align}\label{eq:lagerlof_system} 
    f(G_t,A_t,L_t)\coloneqq
    \begin{cases} 
        G_{t+1} = \big(e(G_t)+\rho\tau\big)\min(\theta L_t,\hat L), \\
        A_{t+1} = (1+G_{t+1})A_t, \\
        L_{t+1} = n(G_t,A_t,L_t)L_t. 
    \end{cases} 
\end{align} 
%
The fertility policy in \eqref{eq:lagerlof_system} is obtained by substituting \eqref{eq:lagerlof_choices}, \eqref{eq:lagerlof_e} into the household solution \eqref{eq:optimal_nt}
%
\begin{equation}\label{eq:lagerlof_optimum} 
    n(G_t,A_t,L_t) = \begin{cases} 
                        \dfrac{\gamma}{\tau+e(G_{t+1})}, & \dfrac{\tilde c}{1-\gamma}\leq z_t, \\[10pt] \dfrac{1-\frac{\tilde c}{z_t}}{\tau+e(G_{t+1})}, & \tilde c\leq z_t<\dfrac{\tilde c}{1-\gamma}, \\[10pt] 0, & z_t<\tilde c. 
                     \end{cases} 
\end{equation}
%
The $\tilde{c}/(1-\gamma)\leq z_t$ and $\tilde{c}\leq z_t < \tilde{c}/(1-\gamma)$ conditions correspond to the household optimum when subsistence consumption is attainable. 
The third condition corresponds to the zero-fertility case, used in \citet{Lagerlof2006}, for states of the economy in which the potential income $z_t$ falls below subsistence $\tilde{c}$. 

As discussed in the previous section, the choice variables are entirely determined by the state variables as 
%
\begin{subequations}\label{eq:lagerlof_observables} 
    \begin{align} 
        e_t &= e(G_t) = \max\!\Big(0,-\rho\tau+\sqrt{\tau(1-\rho)G_t}\Big), \label{eq:lagerlof_education} \\
        z_t &= z(G_t,A_t,L_t) = h(e(G_t),G_t)^{\alpha} \bigg(\frac{A_t}{L_t}\bigg)^{1-\alpha}, \label{eq:lagerlof_income} \\
        h_t &= h(e(G_t),G_t) = \frac{e(G_t)+\rho\tau}{e(G_t)+\rho\tau+G_t}. \label{eq:lagerlof_capital}
    \end{align} 
\end{subequations} 
%
In this work, the GWL model is defined as the non-smooth, nonlinear map f from \eqref{eq:lagerlof_system}, combined with the fertility policy \eqref{eq:lagerlof_optimum} and the choice variables \eqref{eq:lagerlof_observables}.

Figure~\ref{fig:lag_sim} illustrates the trajectory of the GWL model under the calibration of parameters
%
\begin{equation}\label{eq:lagerlof_parameters}
    \tau=0.15,\qquad
    \rho=0.879,\qquad
    \alpha=0.6,\qquad
    \theta=1,\qquad
    \hat L=11.42,\qquad
    \gamma=\tau(1-\rho)(1+\hat L),
\end{equation}
%
and for the initial condition
%
\begin{equation}\label{eq:lagerlof_initial_condition}
    (G_0,A_0,L_0)=(0.048, 0.87, 0.364),
\end{equation}
%
both of which were specified in \citet{Lagerlof2006}.
%
\begin{figure}
        \centering
        \includegraphics[keepaspectratio, width = \textwidth]{figures/Fig1.pdf}
        \caption{\samd{An illustrative trajectory of the GWL model \eqref{eq:lagerlof_system}-\eqref{eq:lagerlof_observables}, replicating the simulation shown in \citet{Lagerlof2006}. Panels (a), (b), and (c) show the evolution of an economy through time (red curves) using its endogenous state variables $G_t$, $A_t$, and $L_t$, respectively. Similarly, panels (d), (e), and (f) show the endogenous choice variables $e_t$, $z_t$, and $h_t$, respectively, which are determined from the state variables at the same time. Panels (a), (c), and (e) also depict the thresholds that determine the relevant branch of \eqref{eq:lagerlof_system} at each time-step. }
        An illustrative trajectory of the GWL model \eqref{eq:lagerlof_system}-\eqref{eq:lagerlof_observables}, replicating the calibration \eqref{eq:lagerlof_parameters}-\eqref{eq:lagerlof_initial_condition} shown in \citet{Lagerlof2006}.
        The economy, represented by the endogenous state vector $X_t=(G_t,A_t,L_t)$ (panels (a)-(c)) and the endogenous variables $e_t$, $z_t$, $h_t$ (panels (d)-(f)), undergoes the three stages predicted by unified growth theory in \citet{Galor2000}.
        A prolonged period of Malthusian stagnation (approximately from generation $t=0$ to $t=40$) is followed by the post-Malthusian transition (from generation $t=41$ to $t=50$) before reaching the sustained growth state (from $t=50$ to $t=60$). 
        Notice that the vertical axis of panel (b), showing the time evolution of the level of technology $A_t$, is in logarithmic scale; the linear path translates in an (unbounded) exponentially growing level of technology.
        The dashed lines in panels (a) and (c) mark the thresholds \eqref{eq:education_threshold} and \eqref{eq:population_threshold}, respectively. 
        The dashed lines in panel (e) (see inset) mark the threshold $\tilde c/(1-\gamma)$ and $\tilde c$ for the branches of the piecewise defined fertility policy \eqref{eq:optimal_nt}.}
        \label{fig:lag_sim}
\end{figure}

\samd{Figure~\ref{fig:lag_sim} reproduces the simulation of \eqref{eq:lagerlof_initial_condition} performed in \cite{Lagerlof2006}, now using the three-dimensional laws of motion \eqref{eq:lagerlof_system}. Panels (a)--(c) depict the evolution of the endogenous state variables of the economy $G_t$, $A_t$, and $L_t$, while panels (d)--(f) show the corresponding choice variables at each time-step. As discussed in \cite{Lagerlof2006}, this trajectory exemplifies the three stages of an economy under unified growth theory. In particular, the economy begins with a period of Malthusian stagnation, which is characterized by slow-moving technological progress $G_t$, a small population in $L_t$, and zero education $e_t$ invested in the future generations. Moreover, the potential income $z_t$ is between $\tilde z$ and $\tilde c$ which constrains the fertility and, therefore, the growth of the population. However, after about 25 generations, the level of technology is sufficiently large to allow $z_t$ to make brief excursions above $\tilde z$. As a result, the economy enters a stage of post-Malthusian growth wherein the population grows faster, which in turn leads to greater technological progress. Around generation 40, the technological progress crosses the threshold \eqref{eq:education_threshold}, meaning that education becomes positive. Accordingly, $A_t$ grows much more quickly, allowing $z_t$ to cross, and remain above $\tilde z$, allowing the economy to reach the maximal fertility branch of \eqref{eq:optimal_nt}. During this period, population exceeds the scale-effect threshold \eqref{eq:population_threshold}, leading to a saturation in both $G_t$ and $L_t$. This transition, marks the final modern-growth stage of the economy which is characterized by a stable population and unbounded improvement in the level of technology.   }

%
\sout{Panels (a)--(c) (left column) display the evolution of the endogenous state vector $X_t=(G_t,A_t,L_t)$ through $60$ generations (red curves).
Notice the presence of the population threshold \eqref{eq:population_threshold} (panel (c), dashed gray line) and the the technological growth threshold \eqref{eq:education_threshold} (panel (a), dashed gray line).
Panels (d)-(f) (right column) display the evolution of the endogenous choice variables implied by this state, $e_t$, $z_t$, and $e_t$ (red curves).
Notice the presence of the $z_t=\tilde c/(1-\gamma)$ and  $z_t=\tilde c$ thresholds (panel (e) and its inset, dashed gray lines) which determine the branch for the piecewise define fertility policy \eqref{eq:lagerlof_optimum}.
The shown trajectory reproduces the three stages emphasized in unified growth theory. 
In the Malthusian trap (approximately from generation $t=0$ to generation $t=40$), technological progress $G_t$ stagnates (panel (a)), population $L_t$ remains small (panel (b)) and the education invested in children $e_t$ is zero (panel (d)).
In the post-Malthusian transition (approximately from generation $t=41$ to generation $t=50$), technological progress is large enough to cross the critical threshold \eqref{eq:education_threshold} (panel (a), red curve crossing the dashed gray line) which, in turn, makes the invested education to become positive (panel (d), lift-off of the red curve from the horizontal axis). 
In the modern-growth regime (from generation $t=51$ to $t=60$), the potential income is larger than the threshold $\tilde c/(1-\gamma)$ (inset of panel (e), red curve crossing the upper dashed gray line after the oscillations) which, in turn, makes the subsistence constraint no longer binding the fertility decision. 
As a result of the erosion of human capital (panel (f)), technological progress (panel (a)), population growth (panel (c)) and education of children (panel (d)) slow down and eventually halt to a steady-state (red curves become and stay flat).}



\subsection{Primer on the dynamics of iterated maps}\label{subsec:dynamical_systems_theory}
What we need so far:
\begin{itemize}
    \item definition of fixed points;
    \item definition of invariant set;
    \item definition of stability of fixed points;
    \item definition of switching manifolds;
    \item brief intro to bifurcation analysis with reference to the literature.
\end{itemize}
%

\section{Dynamics of the GWL model}\label{sec:analysis}
The quantitative exercise in \citet{Lagerlof2006} provides a benchmark transition from Malthusian stagnation to the modern-growth regime of the GWL model \sout{(see Figure~\ref{fig:lag_sim}). }\samb{; see again Figure~\ref{fig:lag_sim}.}
Under Lagerl\"{o}f's calibration of parameters \eqref{eq:lagerlof_parameters} and initial conditions \eqref{eq:lagerlof_initial_condition}, the model demonstrates the typical three-stage progression of an economy under unified growth theory. \sout{the familiar unified growth sequence: a long period of near-subsistence income, slow technological progress, no education, and limited population growth is followed by rising investment in education, faster technological progress, and sustained growth.}

A natural question to ask is how generic this transition is for economies starting at different initial conditions or under different parameters calibration?
In this section we address such question by studying the dynamics of the map $f$ defined in \eqref{eq:lagerlof_system}--\eqref{eq:lagerlof_observables}. 
Specifically, we investigate how the qualitative behavior observed in Figure \ref{fig:lag_sim} changes for economies characterized by different initial growth rates $G_0$, levels of technology $A_0$ and population $L_0$ or under different choice of parameters such as the tradeoff between the consumption and human capital $\gamma\in(0,1)$.

\samd{In this section, we explore whether the trajectory of any economy exhibits the same three stages as shown in Figure~\ref{fig:lag_sim}. In particular, we first investigate the short-term behavior of economies with different initial endogenous state variables $G_0$, $A_0$, and $L_0$. Moreover, we also discuss the different long-term behaviors of economies and particularly highlight the effect of the particular choice of the trade-off parameter $\gamma_c$. In this way, we give a more complete characterization of the different evolutions of economies in the GWL model. }

\subsection{Collapsing economies}\label{subsec:collapse}
We begin our analysis by studying the persistence of the benchmark transition, shown by Lagerl\"{o}f, for economies starting with initial conditions different from those prescribed in \eqref{eq:lagerlof_initial_condition}.
We seek to establish whether the transition from Malthusian stagnation to sustained growth is a global property of the trajectories generated by the GWL model, or rather the result of the choices made in \cite{Lagerlof2006}.
%
\samb{see my comment; important}
\begin{definition}\label{def:population_collapse}
    Let $X_t=(G_t,\,A_t,\,L_t)\in \mathbb{R}^3_{\geq0}$ and $f:\mathbb{R}^3_{\geq0}\to \mathbb{R}^3_{\geq0}$ given by \eqref{eq:lagerlof_system}--\eqref{eq:lagerlof_observables}, then the set
    %
    \begin{equation}\label{eq:population_collapse}
        U_c\coloneqq\{X_t\in \mathbb{R}^3_{\geq0} \;:\;L_t = 0,\,G_t\geq0\;\text{and}\;A_t > 0\}\,,
    \end{equation}
    %
    identifies states of population collapse.
\end{definition}

As a starting point of our study, we repeat the benchmark simulation performed in \cite{Lagerlof2006} using the same parameter values, but now introduce a small perturbation to the initial condition.
In particular, we keep the same initial rate of technological progress $G_0=0.048$ and level of technology $A_0=0.87$ as in \eqref{eq:lagerlof_initial_condition}, but now increase the initial population to $L_0=0.664$. \sout{$L_0=0.364$ by $0.3$\samb{, so that $L_0=0.664$}}.


%
\begin{figure}
        \centering
        \includegraphics[keepaspectratio, width = \textwidth]{figures/greentrajectory.pdf}
        \caption{Time evolution of the GWL model \eqref{eq:lagerlof_system}-\eqref{eq:lagerlof_observables} under the same calibration of parameters \eqref{eq:lagerlof_parameters} (as prescribed in \citet{Lagerlof2006}), but starting at the perturbed initial condition $X_0=(G_0,A_0,L_0)=(0.048, 0.87, 0.664)$ (which preserves Lagerl\"{o}f's prescription \eqref{eq:lagerlof_initial_condition} for the initial level of technological progress $G_0$ and level of technology $A_0$ except with a slightly larger initial population $L_0$). 
        The economy undergoes immediate population collapse $L_1=0$ (panel (c)) after the initial condition.
        As a result, the rate of technological progress $G_t$ (panel (a)) increases slightly after one generation before reducing to zero at the second generation.
        This in turn leaves the level of technology $A_t$ (panel (b)) at a steady-state $A^*=0.95$ until the end of the simulation.}
        \label{fig:lagerlof_perturbation}
\end{figure}
%
Figure~\ref{fig:lagerlof_perturbation} shows the resulting trajectory of the economy starting from the new initial condition $(G_0,A_0,L_0)=(0.048,0.87,0.664)$. \samd{Despite the close proximity to \eqref{eq:lagerlof_initial_condition}, this economy experiences a population collapse after only one generation.}
The trajectory of the economy starting at the perturbed initial condition $X_0=(G_0,A_0,L_0)=(0.048,0.87,0.664)$ (green path) experiences population collapse after one generation. 
This contrasts sharply with the benchmark trajectory in Figure~\ref{fig:lag_sim}, which starts from $X_0=(G_0,A_0,L_0)=(0.048,0.87,0.364)$ (red path) and transitions out of the Malthusian regime toward sustained growth.
Importantly, the trajectory of the perturbed economy (green path) stays in the collapsed state at any $t\geq1$, meaning that once this economy undergoes population collapse it can not recover. Notice how at $t=1$ the population in the economy vanishes and remains in such state for the rest of the simulation, i.e. $L_t=0$ for any $t\geq1$ (panel (c)). From \eqref{eq:lagerlof_system} it is easy to see that if $L_t=0$ then $G_{t+1}=0$ (panel (a)). Consequentially, if $G_{t+1}=0$ then $A_{t+1}=A_t$ (panel (b)).
%
\begin{observation}\label{obs:invariant_set}
    The set of population collapse $U_c$ is invariant under the GWL laws of motion \eqref{eq:lagerlof_system}--\eqref{eq:lagerlof_observables}. That is, once an economy enters $U_c$, it remains in $U_c$. This follows trivially from \eqref{eq:lagerlof_choices} and \eqref{eq:lagerlof_optimum}.
\end{observation}

To characterize the economic mechanism that leads an economy into a state of population collapse we study the evolution of the choice variables $e_t$, $h_t$ and $z_t$. 
From the fertility rule \eqref{eq:lagerlof_optimum}, observe that the optimal number of children $n^*_t$ is set to zero whenever the potential income \samb{$z_t$} falls below the subsistence requirement $\tilde c$. 
Hence, for any economy with population, $L_t>0$, the condition \(z_t<\tilde c \) implies \sout{${n(G_t,A_t,L_t)=0}$} \samb{$n_t^*=0$} and therefore $L_{t+1}=0$.
The simulation in Figure~\ref{fig:lagerlof_perturbation} thus displays a trajectory that belongs to a broader region of initial states that lead to population collapse.
%
\begin{proposition}\label{prop:population_collapse_1}
     Let $X_t\in \mathbb{R}^3_{\geq0}$, $U_c\subset\mathbb{R}^3_{\geq0}$ given by \eqref{eq:population_collapse} and $f:\mathbb{R}^3_{\geq0}\to \mathbb{R}^3_{\geq0}$ given by \eqref{eq:lagerlof_system}--\eqref{eq:lagerlof_observables}, then there exists an open and unbounded subset
     %
     \begin{equation}\label{eq:set_population_collapse_1}
         U_1\coloneqq\{X_t\in\mathbb{R}^3_{\geq0}\;:\; f(X_t)\in U_c\}\,.
     \end{equation}
     %
     of states that undergo population collapse after one generation.
\end{proposition}
%
\samb{I believe it would be beneficial to move Prop. 2 to after the description of Fig4. Furthermore, I would merge Prop. 2 and 3, along with all the asymptotic descriptions, into a single technical proposition that formalizes the visual findings AFTER showing those visual findings. This would drastically improve the flow of this subsection, which is currently interrupted by abstract, technical, and overly detailed results; see my comments for more details.}
\begin{proof}
     From \eqref{eq:lagerlof_optimum}, the subsistence condition $z_t<\tilde{c}$ binds any state of the economy to a population collapse $L_{t+1} = 0$.
     Using \eqref{eq:lagerlof_capital}, we see that this is true for any \sout{$X_t=(G_t,\,A_t,\,L_t)$}\samb{$G_t$, $A_t$ and $L_t$ such that} \sout{that satisfies}
     %
     \begin{equation}\label{eq:population_collapse_1}
         0\leq A_t<L_t(\tilde{c}\,h(e(G_t),G_t)^{-\alpha})^{\frac{1}{{1-\alpha}}}\,.    
     \end{equation}
     %
     By setting 
     %
     \begin{equation*}
         U_1\coloneqq\Big\{X_t\in\mathbb{R}^3_{\geq0}\;:\;A_t<L_t(\tilde{c}\,h(e(G_t),G_t)^{-\alpha})^{\frac{1}{1-\alpha}}\Big\}\,,
     \end{equation*}
     %
     the proof follows.
     \qed
\end{proof}

Proposition~\ref{prop:population_collapse_1} shows that population collapse is not an isolated numerical outcome. 
There is an open and unbounded region $U_1$ of initial states from which the economy undergoes population collapse after one generation, which, by Observation~\ref{obs:invariant_set} is unrecoverable. 
Hence the benchmark transition from Malthusian stagnation to sustained growth is not a global prediction of the GWL model.

Motivated by the existence of $U_1$, we next examine how collapse outcomes vary with the initial level of technology $A_0$ and the initial population $L_0$.
We keep the initial rate of technological progress $G_0=0.048$ fixed at Lagerl\"of's value while varying the remaining two state variables over a grid \((A_0,L_0)\in(0,5]\times(0,5].\)
For each point in this grid, we simulate the trajectory generated by the law of motion $f$ in \eqref{eq:lagerlof_system}--\eqref{eq:lagerlof_observables}.
We then record whether and when the trajectory enters the collapse set $X_c$. 
This procedure allows us to identify regions of initial conditions that lead to population collapse after one or more generations.

Figure~\ref{fig:population_collapse_subspace} displays the resulting classification of initial conditions in the two-dimensional $(L_0,A_0)$-slice, at $G_0=0.048$, of the full state space $\mathbb R^3_{\geq0}$ .
%
\begin{figure}
        \centering
        \includegraphics[keepaspectratio, width = \textwidth]{figures/FIG4.pdf}
        \caption{\samb{To long and repetitive. I would completely rewrite this caption.} Two-dimensional $(A_0,L_0)-$slice of the state space $\mathbb{R}^3_{\geq0}$ at $G_0=0.048$. Depicted is the classification of initial conditions for economies with different starting level of technology $A_0$ and population $L_0$ at a fixed growth rate $G_0=0.048$. \samb{For each point in the grid, we iterate the GWL model \eqref{eq:lagerlof_system}--\eqref{eq:lagerlof_observables} and classify the results by color: blue indicates collapse in one generation, pink in two generations, and white indicates no collapse for up to 60 generations} For each point in the grid we iterate \samb{the GWL model} \eqref{eq:lagerlof_system}--\eqref{eq:lagerlof_observables} forward starting from such condition and record the state of the economy at the end of the simulation (i.e. $T=60$ generations).
        The blue shaded region depicts initial conditions that undergo population collapse after one generation. The blue
        line represents the restriction of the boundary of the set $U_1$ \sout{(see Proposition~\ref{prop:population_collapse_1})}, obtained from \eqref{eq:population_collapse_1}, at $G_0=0.048$.
        The pink shaded region depicts initial conditions that undergo population collapse after two generations. The pink (implicit) curve represents the restriction of the boundary of the set $U_2$ \sout{(see Proposition~\ref{prop:population_collapse_2}), obtained from \eqref{eq:subset_M2_1}-\eqref{eq:subset_M2_6} (see Appendix~\ref{appenix_B}),} \sout{at $G_0=0.0048$.}
        The gray shaded region depicts initial conditions that do no undergo population collapse within the simulated time horizon $T=60$; this set contains the slice of the set $U_{\infty}$ \sout{(see Definition~\ref{def:set_population_survival})} \sout{of population survival at $G_0=0.048$. }
        The trajectory proposed in \cite{Lagerlof2006}, associated to the initial condition \eqref{eq:lagerlof_initial_condition} prescribed by Lagerl\"{o}f, corresponds to the red circle and is located within the gray region. As a result, an economy starting at such red circle will iterate forward without undergoing population collapse, as shown in Figure~\ref{fig:lag_sim}.
        \sout{Notice how such initial condition is close to the boundary of $U_1$ (blue line).}
        Indeed, a relatively larger starting population than the one prescribed by Lagerl\"{o}f in \eqref{eq:lagerlof_initial_condition} is enough to bring the initial condition (green circle) inside the $U_1$ set (blue shaded region).
        The resulting trajectory will thus undergo population collapse within one generation, as shown in Figure~\ref{fig:lagerlof_perturbation}.}
        \label{fig:population_collapse_subspace}
\end{figure}
%
The figure identifies three disjoint regions corresponding to qualitatively different outcomes for the economies' trajectories. 
Initial conditions in the white region do not undergo the collapse set over the simulated horizon of $T=60$ generations.
Initial conditions in the blue region undergo population collapse after one generation. 
This region corresponds to the $(A_0,L_0)-$slice, at $G_0=0.048$, of the set $U_1$ defined in \eqref{eq:set_population_collapse_1}.
Initial conditions in the pink region undergo population collapse after two generations.
The presence of this region reinforces the main implication of Proposition~\ref{prop:population_collapse_1}: that is, under Lagerl\"of's functional forms for human capital and technological progress, the benchmark transition to sustained growth is not a global prediction of the GWL model but a result of the calibrations \eqref{eq:lagerlof_parameters} and \eqref{eq:lagerlof_initial_condition}. 
To illustrate this point, Figure~\ref{fig:population_collapse_subspace} also reports the location of Lagerl\"of's initial condition, $X_0=(G_0,A_0,L_0)=(0.048,0.87,0.364)$, shown by the red circle, and the perturbed initial condition used in Figure~\ref{fig:lagerlof_perturbation}, shown by the green circle. 
The figure shows that, while holding $G_0=0.048$ fixed, sufficiently high starting population $L_0$, relative to the level of technology $A_0$, can lead to population collapse after either one or two generations.
%C
\begin{proposition}\label{prop:population_collapse_2}
    Let $X_t\in \mathbb{R}^3_{\geq0}$, $U_c\in \mathbb{R}^3_{\geq0}$ given by \eqref{eq:population_collapse}, $f:\mathbb{R}^3_{\geq0}\to \mathbb{R}^3_{\geq0}$ given by \eqref{eq:lagerlof_system}--\eqref{eq:lagerlof_observables} and $U_1\subset \mathbb{R}^3_{\geq0}$ given by \eqref{eq:set_population_collapse_1}, then there exists an open subset
    %
    \begin{equation}\label{eq:set_population_collapse_2}
        U_2\coloneqq\{X_t\in \mathbb{R}^3_{\geq0}\;:\; f(X_t)\not\in U_c,\,f^2(X_t)\in U_c\}\,.
    \end{equation}
    %
    that is disjoint from $U_1$.
\end{proposition}
%
\begin{proof}
    See Appendix \ref{appenix_B}.
\end{proof}
%
An immediate result of Proposition \ref{prop:population_collapse_2} is the following.
%
\begin{observation}\label{obs_M2_preimage_M1}
    The image of $U_2$ is contained in $U_1$, i.e. $f(U_2)\subset U_1$.
    To see this notice that $X_2 = f(X_1) = f^2(X_0)$.
    If $X_0\in U_2$ then $X_2 = f(X_1)\in X_c$. 
    From \eqref{eq:set_population_collapse_1} it follows that $X_1=f(X_0)\in U_1$. 
\end{observation}

Observe how in Figure~\ref{fig:population_collapse_subspace}, in the $G_0=0.048$ slice of the state space $\mathbb{R}^3_{\geq0}$, population collapse either occurs within one (blue region) or two generations (pink region) or it does not occur at all (white region).
We next ask whether regions of higher-order population collapse are hidden in other slices of the state space.
To do so, we enlarge the grid for the initial level of technology $A_0$ and population $L_0$ to $(0,250]\times(0,25]$ and we now the initial rate of technological progress $G_0$ vary over the range $(0,0.5]$. 
Specifically, for each selected value of $G_0\in(0,0.5]$, we obtain a two-dimensional slice $(A_0,L_0)\in(0,250]\times(0,25]$ of the state space $\mathbb{R}^3_{\geq0}$ and repeat the simulation procedure used to construct Figure~\ref{fig:population_collapse_subspace}. 

Figure~\ref{fig:subspace_slicing} reports the results of this extended numerical exercise across different values of $G_0$.
%
\begin{figure}
        \centering
        \includegraphics[keepaspectratio, width = \textwidth]{figures/Fig5_slices.pdf}
        \caption{Two-dimensional $(A_0,L_0)-$slices of the state space $\mathbb{R}^3_{\geq0}$ at different values of the growth rate: (a) $G_0=0.048$ (this is the same slice shown in Figure~\ref{fig:population_collapse_subspace} but for a larger region); (b) $G_0=0.35$; (c) $G_0=0.42$ and; (d) $G_0=0.5$.
        Notice how, at small values of $G_0$ (panel (a)) the boundary of the set $U_1$ of population collapse after one iteration (blue curve) decreases its slope (see Eq.~\eqref{eq:asymp_at_0}); conversely, at large values of $G_0$ (panel (d)) the slope increases (see Eq.~\eqref{eq:asymp_at_infty}).
        On the other hand, the set $U_2$ of population collapse after two iterations (pink shaded regions) is large for small values of $G_0$ (panel (a)) and shrinks as $G_0$ increases (panels (b) and (c)) eventually disappearing from the $(A_0,L_0)-$slice at larger values of $G_0$ (panel (d)).}
        \label{fig:subspace_slicing}
\end{figure}
%
Notice how the collapse regions $U_1$ and $U_2$ change as the initial rate of technological progress $G_0$ varies in $(0,0.5]$. 
Three qualitative patterns emerge.
First, the set of population collapse after one generation $U_1$ (blue region) expands as $G_0$ increases from Lagerl\"of's benchmark value, $G_0=0.048$, to $G_0=0.5$.
This easily follows from \eqref{eq:population_collapse_1} for the boundary of the $U_1$ given by
%
\begin{equation}\label{eq:boundary_M1}
    A_0 = L_0\left(\tilde{c}h(e(G_0),G_0)^{-\alpha}\right)^{\frac{1}{1-\alpha}}.
\end{equation}
%
Thus, at any slice $G_0$, the boundary of $U_1$ gives a linear relationship between starting level of technology $A_0$ in an economy and its population $L_0$.
Specifically, from \eqref{eq:lagerlof_education} and \eqref{eq:lagerlof_capital}, we get
%
\begin{align*}
    e(G_0)\to 0,\qquad h(G_0)\to 1
    \quad &\text{as}\quad G_0\to0, \\
    e(G_0)\to\infty,\qquad h(G_0)\to0
    \quad &\text{as}\quad G_0\to\infty.
\end{align*}
%
It follows that
%
\begin{subequations}\label{eq:asymptotic_equalities}
\begin{align}
     \left(\tilde{c}h(G_0)^{-\alpha}\right)^{\frac{1}{1-\alpha}}
     \to
     \tilde{c}^{\frac{1}{1-\alpha}}
     \quad &\text{as}\quad G_0\to0, \label{eq:asymp_at_0}\\
     \left(\tilde{c}h(G_0)^{-\alpha}\right)^{\frac{1}{1-\alpha}}
     \to
     \infty
     \quad &\text{as}\quad G_0\to\infty. \label{eq:asymp_at_infty}
\end{align}
\end{subequations}
%
Hence, at low values of $G_0$, the boundary of the region of population collapse $U_1$, after one generation, approaches $A_0=\tilde{c}^{\frac{1}{1-\alpha}}L_0$.
As $G_0$ increases, the boundary rotates upward in the $(L_0,A_0)$-plane. Consequently, over any fixed bounded plotting window, a larger share of initial conditions lies below the boundary and therefore belongs to the one-generation collapse region.

Second, Figure~\ref{fig:subspace_slicing} shows that the set of population collapse after two generations $U_2$ (pink region), shrinks as $G_0$ increases, eventually vanishing from the state space. 
In contrast to the boundary of $U_1$, the boundary of $U_2$ is nonlinear and implicit (see Appendix~\ref{appenix_B}). Hence there is no simple boundary expression analogous to \eqref{eq:boundary_M1}. 
Nevertheless, the pattern has a clear economic interpretation: when the initial rate of technological progress is sufficiently high, economies that do not collapse after one generation are less likely to enter the collapse region $U_1$ in the following period.

Third, we do not observe regions in which collapse occurs after three or more generations. In the simulated subset $ (G_0,A_0,L_0)\in(0,0.5]\times(0,250]\times(0,25]$ the only collapse regions identified are $U_1$ (blue), corresponding to collapse after one generation, and $U_2$ (pink), corresponding to collapse after two generations. 
Initial conditions in the white region do not enter the collapse set over the simulated horizon. This motivates the following definition.
%
\begin{definition}\label{def:set_population_survival}
     Let $X_t\in \mathbb{R}^3_{\geq0}$, $U_c\in \mathbb{R}^3_{\geq0}$ given by \eqref{eq:population_collapse} and $f:\mathbb{R}^3_{\geq0}\to \mathbb{R}^3_{\geq0}$ given by \eqref{eq:lagerlof_system}--\eqref{eq:lagerlof_observables}, then the set
     %
     \begin{equation}\label{eq:set_population_survival}
         U_{\infty} = \{X_t\in \mathbb{R}^3_{\geq0} \;:\; f^{n}(X_t)\neq U_c\,\;\forall n\in\mathbb{Z}\}\,,
     \end{equation}
     %
     identifies states that do not collapse.
\end{definition}
%
The set $U_{\infty}$ identifies the white regions in Figures~\ref{fig:population_collapse_subspace} and \ref{fig:subspace_slicing}. 
To characterize the behavior of economies that belong to such set, we now turn to the long-run dynamics of the GWL model.

Proposition~\ref{prop:section_3_1} formalizes our observations.
\begin{proposition}     \label{prop:section_3_1}

\end{proposition}
%
\begin{proof}
    See Appendix \ref{appenix_B}.
\end{proof}
%
\subsection{Steady-states of the economy}\label{subsec:fixed_points}
\sout{We start by investigating the existence of fixed points of the GWL model in the full state space $\mathbb{R}^3_{\geq0}$. }\samb{In this section, we investigate the existence and stability of the fixed points of the GWL model.}
\sout{Recall that, by Definition \textbf{[INSERT REF.]}, a fixed point $X^*=(G^*,A^*,L^*)\in\mathbb{R^3}_{\geq0}$ of the map $f$ must satisfy the identity $X^*=f(X^*)$.}
\samb{In particular, w}\sout{W}e focus on economically meaningful states with $A^*>0$.

Notice that, from the technology law of motion in \eqref{eq:lagerlof_system}, $A_{t+1}=A_{t}$ if and only if $ G_{t+1}=0$.
Thus, at any fixed point with $A^*>0$, the condition $A_{t+1}=A_t$ implies $G_{t+1}=0$. 
Since a fixed point also requires $G_{t+1}=G_t$, any fixed point must also satisfy $G^*=0$.
Under Lagerl\"of's specification for technological progress \eqref{eq:lagerlof_choices}, $e(G_t)\geq0$ and $\rho\tau>0$, imply that $e(G_t)+\rho\tau$ is strictly positive. 
Therefore $G_{t+1}=0$ if and only if $L_t=0$.
We conclude that any fixed point $X^*\in \mathbb{R}^3_{\geq0}$ of the map $f$ given in \eqref{eq:lagerlof_system}-\eqref{eq:lagerlof_observables} is characterized by
%
\begin{equation}\label{eq:line_fixed_points}
     U_c\ni X^*=(G^*,A^*,L^*)=(0,A^*,0)\,,\quad A^*>0.
\end{equation}
%
\begin{observation}\label{obs:malthusian_trap}
     The Malthusian regime (slow technological progress $G_t\gtrapprox0$ and small population level $L\gtrapprox0$) is not a fixed point of Lagerl\"{o}f's dynamical system.
\end{observation}
%
This observation clarifies the role of the Malthusian regime in the benchmark simulation proposed in \cite{Lagerlof2006}.
While an economy might spend several generations in a regime of slow growth and low education, such regime is not a steady state of the GW model. 
This means that no economy can ever remain in a regime of Malthusian stagnation.
The only states that remain exactly unchanged over time are those of population collapse $U_c$ (see Definition~\ref{def:population_collapse}).
A natural question to ask is whether the fixed points $X^*\in U_c$ attract nearby states of the economy.
Specifically, we seek to establish wether economies that are close to a fixed point $X^*$ will inevitably undergo population collapse or they escape to alternative regimes (e.g. the sustained growth state showed in Figure \ref{fig:lag_sim}).

To do so we study the local stability of the fixed points $X^*$ given by \eqref{eq:line_fixed_points}.
\sout{By Definition \textbf{[INSERT REF.]},}\samb{As detailed in Section~\ref{subsec:dynamical_systems_theory},} this is achieved by studying the eigenvalues of the linearization of the map \eqref{eq:lagerlof_system}-\eqref{eq:lagerlof_observables} evaluated at the fixed points, i.e. $D_f(X^*)$.
Notice that, by virtue of $f$ having multiple (switching) branches, we must first determine which of those need to be selected to study local (linear) stability of the fixed points $X^*$.
From \eqref{eq:line_fixed_points}, the set of states in $\mathbb{R}^3_{\geq0}$ that are close to $X^*$ are characterized by small (positive) values of growth $G_t\gtrsim0$ and population $L_t\gtrsim0$.

\greta{First, $G_t$ lies below the education threshold, so $e(G_t)=0$. Second, $L_t$ lies below the population scale-effect threshold, so $\min\{\theta L_t,\hat L\}=\theta L_t$. Third, since $A_t/L_t\to\infty$ as $L_t\to0^+$, potential income satisfies $z(G_t,A_t,L_t)\to\infty$; hence the subsistence constraint is locally slack and the high-income fertility branch applies. 
Finally, for $L_t$ sufficiently small, $G_{t+1}=\rho\tau\theta L_t$ also lies below the education threshold, so $e(G_{t+1})=0$.}
\greta{In particular, we can define such set to be}
%
\greta{
\begin{equation*}
    U^*\coloneqq \bigg\{X_t=(G_t,A_t ,L_t)\in M\;:\;0<G_t<\frac{\rho^2\tau}{1-\rho}\,,\;0<L_t<\frac{\hat{L}}{\theta}\bigg\}\,.
\end{equation*}
}
%
From \eqref{eq:lagerlof_system}-\eqref{eq:lagerlof_observables} it thus follow that the local restriction of $f$ to $U^*$ is
%
\begin{equation}\label{eq:lagerlof_restricted}
    f^*(G_t,A_t,L_t)\coloneqq f|_{U*}(G_t,A_t,L_t)=
     \begin{cases} 
          \rho\tau\theta\,L_t, \\[6pt] 
          (1 + \rho\tau\theta\,L_t)A_t, \\[6pt] 
          \frac{\gamma}{\tau}L_t. 
     \end{cases} 
\end{equation}
%
With this we are able to formalize the description of the fixed points given in \eqref{eq:line_fixed_points}.
%
\begin{proposition}\label{prop:line_of_saddles}
    Let $X_t=(G_t,\,A_t,\,L_t)\in\mathbb{R}^3_{\geq0}$, $U_c\in\mathbb{R}^3_{\geq0}$ given by \eqref{eq:population_collapse} and $f:\mathbb{R}^3_{\geq0}\to\mathbb{R}^3_{\geq0}$ given by \eqref{eq:lagerlof_system}--\eqref{eq:lagerlof_observables}, then the subset
    %
    \begin{equation}\label{eq:line_of_saddles}
        X^*=\{X_t\in U_c\;:\;G_t=0,\,L_t=0\;\text{and}\;A_t>0\}\,,
    \end{equation}
    %
    is a set of saddle points for $f$.
\end{proposition}
%
\begin{proof}
    From direct calculation, the Jacobian of the local restriction of $f$, given in \eqref{eq:lagerlof_restricted}, and evaluated at the fixed points $X^*$ given in \eqref{eq:line_fixed_points}, is
    %
    \begin{equation*}
        D_{f^*}(X^*)=\begin{pmatrix}
                        0 & 0 & \rho\tau\theta\\
                        0 & 1 & \rho\tau\theta\,A^* \\
                        0 & 0 & \gamma/\tau
                     \end{pmatrix}\,,
    \end{equation*}
    %
    whose (real) eigenvalues are $\lambda\in\{0,1,\gamma/\tau\}$. 
    By Definition \textbf{[INSERT REF.]}, any fixed point in $X^*$ corresponds to a saddle point with a stable ($\lambda=0$), center ($\lambda=1$) and unstable ($\lambda=\gamma/\tau>1$) eigenspaces.
    \qed
\end{proof}
\samb{We need to conclude somethign from this proposition: "Proposition~\ref{prop:line_of_saddles} demonstrates the the fixed point $X^*$ is not attracting; that is, ..." }
\subsection{Asymptotic reduction in the modern-growth regime}\label{subsec:modern_growth_regime}
%
The previous section showed that the only steady states of the full three-dimensional GWL model are states of population collapse (see Proposition \ref{prop:line_of_saddles}). 
This raises an apparent tension with the benchmark simulation in Figure~\ref{fig:lag_sim}. 
Along that trajectory, the economy escapes the Malthusian regime and enters a modern-growth regime in which technological growth $G_t$ and population $L_t$ settle on constant values \sout{(Figure \ref{fig:lag_sim}(a-b))}\samb{as shown in panels~(a) and (c) of Figure \ref{fig:lag_sim}}.
However, this state of sustained growth is not a fixed point in the full, three-dimensional state space $\mathbb{R}^3_{\geq0}$, since the level of technology $A_t$ never settles on a fixed value $A^*$\sout{ (see Figure \ref{fig:lag_sim}(c))}\samb{; see Figure \ref{fig:lag_sim}(b)}; if it did, then the state of the economy would result in a population collapse as in \eqref{eq:line_fixed_points}.
In this section we provide an analytical explanation of this observation. 

To start let us observe that, from \eqref{eq:galor_family}, $A_t$ increases monotonically for any $G_{t+1}>0$. This entails that, for any economy whose trajectory satisfies $G_{t+1}>0$, the level of technology $A_t$ will grow unboundedly. 
In this case, an economy will eventually reach a state in which the level of technology $A_t$ becomes sufficiently large so that the subsistence constraint no longer binds.
Along such trajectory, the potential income eventually exceeds the subsistence threshold, so that the state switches to the high-income fertility branch $z_t>\frac{\tilde{c}}{1-\gamma}$ of \eqref{eq:lagerlof_optimum} (see Figure \ref{fig:lag_sim}(f)).

To derive this reduced system, consider trajectories for which the level of technology becomes sufficiently large relative to population.  In this region, the dynamics of $G_t$ and $L_t$ no longer depend on the level of technology $A_t$. The level of technology continues to evolve according to $A_{t+1}=(1+G_{t+1})A_t$, but it is no longer needed to determine the laws of motion for technological growth and population.
In this regime, the system prescribed by Lagerl\"{o}f \eqref{eq:lagerlof_system} has the form
%
\begin{align*}
    f(G_t,A_t,L_t)=
    \begin{cases}
        G_{t+1} &= (e(G_t) + \rho\tau)\min(\theta L_t, \hat L)\,, \\
        A_{t+1} &= (1 + G_{t+1})A_t\,, \\
        L_{t+1} &= \frac{\gamma}{\tau + e(G_{t+1})} L_t\,,
    \end{cases}
\end{align*}
%
where $e(G_t)$ is given by \eqref{eq:lagerlof_education}.
Importantly, in this asymptotic formulation, Lagerl\"{o}f's laws of motion for growth $G_{t+1}$ and population $L_{t+1}$ no longer depend on the level of technology $A_t$. 
As a consequence, one can study the evolution of the state variables $G_t$ and $L_t$ independently from $A_t$. 
In other words, in the asymptotic regime of large enough $A_t$, the income $z_t$ of the agents is above subsistence $\frac{\tilde{c}}{1-\gamma}$ and the state of the economy is entirely described by its growth $G_t>0$ and population $L_t>0$.

The asymptotic dynamics is therefore described by the two-dimensional reduction of the following form
%
\begin{equation}\label{eq:asymptotic_reduction}
f_\infty(G_t,L_t)
    \coloneqq
    \begin{cases}
        G_{t+1} = (e(G_t)+\rho\tau)\min(\theta L_t,\hat L),\\[4pt]
        L_{t+1} = \dfrac{\gamma}{\tau+e(G_{t+1})}L_t,
    \end{cases}
\end{equation}
%
where $e(G_t)$ is given by \eqref{eq:lagerlof_education}.
In this two-dimensional asymptotic reduction \eqref{eq:asymptotic_reduction}-\eqref{eq:lagerlof_education} the system remains piecewise defined. 
In particular, the education policy still identifies the threshold $G_t=\rho^2\tau/(1-\rho)$ given in \eqref{eq:education_threshold}, while the technological-growth identifies the population scale-effect threshold $L_t=\hat L/\theta$ given in \eqref{eq:population_threshold}. 
Furthermore, from \eqref{eq:asymptotic_reduction} observe that the term $e(G_{t+1})$ in $L_{t+1}$ induces an additional, more complicated threshold for both $G_t$ and $L_t$.
%
\begin{figure}
    \centering
    \includegraphics[keepaspectratio, width = \textwidth]{figures/phaseportraitone.pdf}
    \caption{Partition of the state space $\mathbb{R}^2_{>0}$ of the asymptotic reduction $f_{\infty}$ \eqref{eq:asymptotic_reduction}-\eqref{eq:lagerlof_education} of Lagerl\"{o}f's system $f$ \eqref{eq:lagerlof_system}-\eqref{eq:lagerlof_observables}. The six (disjoint) regions $\mathcal A$--$\mathcal F$ are bounded by the switching manifolds $M_{1}$ (green line), $M_2$ (blue line) and $M_3$ (purple curve) given in \eqref{eq:switching_manifolds}.
    The manifolds $M_1$ and $M_2$ intersect at the state $(G_t,L_t)=(\rho^2\tau/(1-\rho),\hat{L}/\theta)$ (green square) whereas $M_2$ and $M_3$ intersect at the state $(G_t,L_t)=(\rho^2\tau/(1-\rho),\rho/\theta(1-\rho))$ (purple square), which is also the point at which the two branches of $M_3$ meet (see Eq.~\eqref{eq:switching_manifold_3}).
    Notice how region $\mathcal B$ is bounded above by $M_1$, below by $M_3$, on the right by $M_2$ and on the left by the $L_t-$axis.
    Similarly, region $\mathcal{C}$ is bounded above by $M_3$, below by the $G_t-$axis, on the right by $M_2$ and on the left by the $L_t-$axis.
    Conversely, regions $\mathcal A$, $\mathcal D$, $\mathcal E$ and $\mathcal F$ are unbounded. \samb{Shall we add $X^*$?}}
    \label{fig:asym_phase}
\end{figure}
%

These threshold boundaries partition the state space $\mathbb R^2_{>0}$ of the asymptotic system into regions with distinct local laws of motion and, therefore, distinct qualitative dynamics.
%
\begin{proposition}\label{prop:switching_manifolds}
    Let $X_t=(G_t,L_t)\in\mathbb{R}^2_{>0}$ and $f_{\infty}$ given by \eqref{eq:asymptotic_reduction}-\eqref{eq:lagerlof_education}, then
    %
    \begin{subequations}\label{eq:switching_manifolds}
        \begin{align}
            M_1 &\coloneqq \bigg\{X_t\in\mathbb{R}^2_{>0}\;:\;L_t=\frac{\hat{L}}{\theta}\bigg\}, \label{eq:switching_manifold_1} \\
            M_2 &\coloneqq \bigg\{X_t\in\mathbb{R}^2_{>0}\;:\;G_t=\frac{\rho^2\tau}{1-\rho}\bigg\}, \label{eq:switching_manifold_2} \\
            M_3 &\coloneqq \Bigg\{X_t\in\mathbb{R}^2_{>0}\;:\;L_t=\min \bigg\{\frac{\rho}{(1-\rho)\theta},\frac{\rho^2\tau}{(1-\rho)\theta\sqrt{\tau(1-\rho)\,G_t}}\bigg\}\Bigg\}, \label{eq:switching_manifold_3}
        \end{align}
    \end{subequations}
    %
    are switching manifolds of $f_{\infty}$.
\end{proposition}
%
\begin{proof}
    The derivation of $M_1$, $M_2$ and $M_3$ can be found by following the steps in Appendix \ref{appenix_B} until \eqref{eq:sub_partition_2_1} by noticing that the partitioned $(G_t,L_t)-$plane in the derivation of \eqref{eq:set_population_collapse_2} and the state space $\mathbb{R}^2_{>0}$ of the reduced $f_{\infty}$ are the same object. 
    Indeed the sets $M_{j}$, $j=1,2,3$ given in \eqref{eq:switching_manifolds} are the boundaries of the sets $\Omega_j$, $j=1,\dots,6$ given in \eqref{eq:partition_omega}. \qed
\end{proof}

\samb{Lets find a name for the intersection points $M_1\cap M_2$ and $M_2\cap M_3$ and lets introduce those in the Proposition. }

The switching manifolds $M_{1,2,3}$, (see Definition \textbf{[INSERT REF.]}) organize the dynamics of the reduced system $f_{\infty}$ by partitioning the state space $\mathbb R^2_{>0}$ into size regions, denoted by $\mathcal A$--$\mathcal F$, with distinct local laws of motion. 
Figure~\ref{fig:asym_phase} illustrates this partition \samb{in the $(G,L)$-plane}.
Notice that, while the switching manifolds $M_1$ (green line) and $M_2$ (blue line) provide trivial thresholds in the $L_t$ and $G_t$ coordinates respectively (see Eqs.~\eqref{eq:switching_manifold_1}-\eqref{eq:switching_manifold_2}), the switching manifold $M_3$ (red curve) is given by the piecewise-continuous curve $L_t=L(G_t)$ whose branches meet at the point $(G_t,L_t)=(\rho^2\tau/(1-\rho),\rho/\theta(1-\rho))$ of the plane (see Eq.~\eqref{eq:switching_manifold_3}).
Furthermore, from Figure~\ref{fig:asym_phase} it is easy to see how only the regions $\mathcal B$ and $\mathcal{C}$ are bounded subsets of $\mathbb{R}^2_{>0}$, as opposed to the remaining ones which are unbounded in either the $G_t$ coordinate (regions $\mathcal E$ and $\mathcal F$), the $L_t$ coordinate (region $\mathcal A$) or both (region $\mathcal D$).

To study the steady-states, depicted in Figure~\ref{fig:lag_sim}(a-b), for the modern-growth regime, we observe that the asymptotic (reduced) dynamics $f_{\infty}$ has different forms depending on which of these six regions the state of the economy $X_t=(G_t,L_t)$ is (see Appendix \ref{appenix_C1}).
By analyzing the reduced dynamics defined for each of the regions $\mathcal A$--$\mathcal F$, we conclude that states $X_t\in\mathbb{R}^2_{>0}$ can either: transition from one region to another or; remain in the same region.

\samb{We should be cautious regarding the boundaries and what they mean $\Gamma$}
This dynamical structure induces a directed graph $\Gamma$, whose nodes represent the six regions $\mathcal A$--$\mathcal F$ partitioning the state space $\mathbb{R}^2_{>0}$ and whose (directed) edges represent the possible transitions among $\mathcal A$--$\mathcal F$ dictated by $f_{\infty}$ (see Appendix \ref{appenix_C2}).
This \emph{transition graph} $\Gamma$, illustrated in Figure~\ref{fig:graph}, summarizes the admissible transitions among regions of the asymptotic reduction of the GWL model.
%
\begin{figure}
    \centering
    \includegraphics[keepaspectratio, scale = 0.5]{figures/samd/transition_graph.pdf}
    \caption{Transition graph $\Gamma$ induced by the asymptotic dynamics $f_{\infty}$. Each node represents one of the six regions $\mathcal A$--$\mathcal F$ that partition the state space $\mathbb{R}^2_{>0}$ through the switching manifolds $M_{1,2,3}$ given by \eqref{eq:switching_manifolds} (see Figure~\ref{fig:asym_phase}).
    The (directed) edges represent each possible transitions of states $X_t$ following one iteration of \eqref{eq:asymptotic_reduction}-\eqref{eq:lagerlof_education}.
    All regions allow for multiple, possible transitions (e.g. states in $\mathcal B$ can either transition into $\mathcal D$ or $\mathcal E$) with the exception of $\mathcal A$ which can only transition into $\mathcal D$.
    Region $\mathcal A$ is also the only one that lacks an incoming edge.
    The pair $\mathcal D$--$\mathcal E$ is the only set of regions that allow transitions in both directions.
    Notice that transitions are allowed only for contiguous regions of $\mathbb{R}^2_{>0}$ with the exception of the transition from $\mathcal D$ to $\mathcal F$.
    The returning transitions for regions $\mathcal C$, $\mathcal D$ and $\mathcal E$ identify the necessary conditions for fixed points of $f_{\infty}$ to exist in $\mathbb{R}^2_{>0}$.}
    \label{fig:graph}
\end{figure}
%
Nodes in $\Gamma$ correspond, one-to-one, to regions $\mathcal A$--$\mathcal F$ of $\mathbb{R}^2_{>0}$ (see Figure~\ref{fig:asym_phase}).
A directed edge connecting two nodes in $\Gamma$ identifies admissible transitions to and from regions $\mathcal A$--$\mathcal F$ of $\mathbb{R}^2_{>0}$ within one generation.
The orientation of the edges indicate the direction with which such transitions occur.
For example nodes $\mathcal B$ and $\mathcal E$ are connected by an edge, whereas $\mathcal B$ and $\mathcal A$ are not connected by an edge.
This means that economies in region $\mathcal A$ cannot transition into economies in region in $\mathcal A$ in one step.
Furthermore, the edge connecting $\mathcal B$ and $\mathcal E$ is oriented from $\mathcal B$ to $\mathcal E$.
This means that economies that are in $\mathcal B$ can transition into $\mathcal E$ in one step, but not viceversa.
An edge with a double orientation entails that economies between those states connected by said edge can transition  in both directions.
This only occurs for the edge connecting regions $\mathcal D$ and $\mathcal E$.

Notice that some nodes in $\Gamma$ have more more than one edge connected to them.
For example, $\mathcal B$ has two outgoing edges connecting it to $\mathcal D$ and $\mathcal E$ and two incoming edges connecting it to $\mathcal C$ and $\mathcal F$.
This means that economies $X_t$, that at time $t$ are in $\mathcal B$, are the evolution, after one generation, of economies $X_{t-1}$ that were either in $\mathcal C$ or $\mathcal F$, at generation $t-1$.
Specularly, these economies will evolve, in the next generation $t+1$, in states $X_{t+1}$ that will either be in $\mathcal D$ or $\mathcal E$.
Node $\mathcal A$ only has one outgoing edge connecting it to $\mathcal D$.
Because $\mathcal A$ has no incoming edge, economies in $\mathcal A$ can not persist in such region and will be in $\mathcal D$ after one generation.

Finally, observe that some edges in $\Gamma$ connect a node to itself.
This is the case for $\mathcal C$, $\mathcal D$ and $\mathcal E$.
This entails that fixed points of the asymptotic reduction $f_{\infty}$ can only exist in regions $\mathcal{C}$, $\mathcal{D}$ and $\mathcal{E}$ of the state space $\mathbb{R}^2_{>0}$.
We remark that the presence of these \emph{returning edges} only give necessary, but not sufficient, conditions for the existence of steady-states.
We focus here on one particular steady-state of the economy in the modern-growth regime, which is the one reported in \cite{Lagerlof2006} (see Figure~\ref{fig:lag_sim}(a-b)).
%

\samb{I think this section would gain from having a subsection 3.3.1 starting here. It could be titled: "Bifurcation analysis of the reduced system" or something.}

\samb{As a shorthand notation, we should introduce the critical $\gamma$-value 
\begin{equation}
\label{eq:gammac}
    \gamma_c=:\tau(1-\rho)(1+\hat{L}).
\end{equation}
=
}


\begin{proposition}\label{prop:asymp_fixed_points}
    Let $f_{\infty}$ given by \eqref{eq:asymptotic_reduction}-\eqref{eq:lagerlof_education}, $\tau$, $\rho$ and $\theta$ given by \eqref{eq:lagerlof_parameters} and $\gamma<\gamma_c$, then
    %
    \begin{equation}\label{eq:fixed_points_asymp}
        X^*=(G^*,L^*)=\bigg(\frac{(\gamma-\tau(1-\rho))^2}{\tau(1 - \rho)}, \frac{\gamma-\tau(1-\rho)}{\theta\,\tau(1 - \rho)}\bigg)\,,
    \end{equation}
    %
    is a stable fixed point of $f_{\infty}$ in $\mathbb{R}^2_{>0}$.
\end{proposition}
%
\begin{proof}
    To see this, one can simply evaluate $f_{\infty}$ at the proposed state $X^*$ given by \eqref{eq:fixed_points_asymp} and observe that $X^*=f_{\infty}(X^*)$, which, by Definition \textbf{INSERT REF.}, gives a fixed point of $f_{\infty}$.
    To determine the stability of $X^*$ one computes the eigenvalues of the Jacobian of $f_{\infty}$, also evaluated at $X^*$. These are
    \begin{equation*}
        -1<\lambda_{1,2}=\frac{\beta}{2}\pm\frac{\sqrt{\beta^2-2}}{2}<1\,,\quad\beta\coloneqq\frac{2\gamma+\tau(1-\rho)}{2\gamma}\,,
    \end{equation*}
    where we omitted the calculations, as they follow the exact same steps as in the proof of Proposition \ref{prop:line_of_saddles}.
    \qed
\end{proof}

The partition $\mathcal A$--$\mathcal F$, and the associated transition graph $\Gamma$, are parameter dependent. 
This dependence follows directly from the construction of the switching manifolds in \eqref{eq:switching_manifolds} (see Proposition \ref{prop:switching_manifolds}).
Some of these parameters come from the GW families of models, specifically the labor share $\alpha\in(0,1)$, the preference weight on children $\gamma>0$, the fixed time cost of child rearing $\tau\geq0$, and the subsistence requirement $\tilde c\geq0$. 
The remaining ones come from the specific GWL model, namely $\rho\in(0,1)$, $\theta>0$, and $\hat L>0$. 
Indeed Figure~\ref{fig:asym_phase} and \ref{fig:graph} are a result of Lagerl\"{o}f's choices for the parameters \eqref{eq:lagerlof_parameters}.
Changing the value of those parameters, even by small perturbations, could change the structure of the partitioning (and hence the graph $\Gamma$ that describes the dynamics $f_{\infty}$) in the modern-growth regime.
Perhaps more importantly, one can clearly see from \eqref{eq:fixed_points_asymp} how the steady-state shown by Lagerl\"{o}f (see Figure~\ref{fig:lag_sim}(a-b)) depend on specific values of $\gamma$, $\tau$, $\rho$ and $\theta$.
In particular, notice that by varying the preference weight $\gamma>0$, we obtain qualitatively different steady-states in the modern-growth regime, as shown in Figure~\ref{fig:different_steady_states}.
%
\begin{figure}
        \centering
        \includegraphics[keepaspectratio, width = \textwidth]{figures/varyinggamma.pdf}
        \caption{Time evolution of the population component $L_t$ (red curves) of an economy $X_t=(G_t,L_t)\in\mathbb{R}^2_{>0}$, in the asymptotic reduction \eqref{eq:asymptotic_reduction}-\eqref{eq:lagerlof_education}, at varying values of the preference weight: (a) $\gamma=0.21<\gamma_c$; (b) $\gamma=\gamma_c$ (this corresponds to Lagerl\"{o}f's calibration \eqref{eq:lagerlof_parameters}); (c) $\gamma=0.24>\gamma_c$.
        In each of the three cases the initial condition for the simulation is kept fixed, and in particular $(G_0,L_0)=(1.0752,9.9109)\in\mathcal E$, where $\mathcal E$ is one of the six regions that partition $\mathbb{R}^2_{>0}$ (see Figure~\ref{fig:asym_phase}).
        Panel (a): at $\gamma=0.21<\gamma_c$ the steady-state population component $L^*$ (blue solid line) of the fixed point $X^*$ (see Proposition~\ref{prop:asymp_fixed_points}) is smaller than the threshold $\hat{L}/\theta$ (green dashed line) set by the switching manifold $M_1$ (see Proposition~\ref{prop:switching_manifolds}); the population trajectory undergoes dampened oscillations as it relaxes onto the steady-state $L^*$.
        Panel (b): at Lagerl\"{o}f's calibration $\gamma=\gamma_c$ the fixed point component $L^*$ and the critical threshold $\hat L/\theta$ coincide (blue solid line) and the population monotonically approaches a different steady-state.
        Panel (c): at $\gamma = 0.24 > \gamma_c$, $L^*$ is larger than $\hat L/\theta$ and the population grows unboundedly (inset (c2)).
        }
        \label{fig:different_steady_states}
\end{figure}

We identified three regions of $\gamma$ values yielding these qualitatively different outcomes.
To discuss these, we keep the initial condition $X_0=(G_0,L_0)=(1.0752,9.9109)$ fixed as we vary the preference weight $\gamma$; notice that, for any $\gamma\in(0,1)$, our choice of initial condition $X_0$ lies in the subregion $\mathcal E\subset\mathbb{R}^2_{>0}$.
For values of the preference weight smaller than the one prescribed by Lagerl\"{o}f $\gamma=\gamma_c$ in \eqref{eq:lagerlof_parameters} (e.g. $\gamma=0.21$ as in Figure~\ref{fig:different_steady_states}(a)) the population $L_t$ (red curve) undergoes damped oscillations as it approaches a steady-state value $L^*$ (blue line) as given in \eqref{eq:fixed_points_asymp}.
At Lagerl\"{o}f's calibration $\gamma=\gamma_c$ (Figure~\ref{fig:different_steady_states}(b)) the population $L_t$ approaches a different steady-state value instead and it does so monotonically.
For values of $\gamma$ larger than Lagerl\"{o}f calibration ((e.g. $\gamma=0.24$ as in Figure~\ref{fig:different_steady_states}(c))) the population $L_t$ never settles onto a steady-state and grows unboundedly (see inset (c2)).
As a result, it is easy to see how the trajectory shown in \citet{Lagerlof2006} is not robust; a small perturbation of $\gamma$ produces qualitatively different economic trajectories in the modern-growth regime.
In particular, we emphasize how economies that reach the modern-growth regime, with a slightly larger preference weight placed on children $\gamma$ than the prescribed value $\gamma_c$, do not settle on a sustained growth steady-state and their population $L_t$ grows without bounds.

%
\begin{figure}
    \centering
    \includegraphics[keepaspectratio, width = \textwidth]{figures/bifdiagram_paper.pdf}
    \caption{Bifurcation diagram (panel (a)) and state space $\mathbb{R}^2_{>0}$ (panels (b1)-(b3)) of the asymptotic reduction $f_{\infty}$ given by \eqref{eq:asymptotic_reduction}-\eqref{eq:lagerlof_education}.
    Panel (a): depicted is the (linear) dependence of the population component $L^*$ (blue line) of the fixed point $X^*$ (see Eq.~\eqref{eq:fixed_points_asymp}) from the preference weight $\gamma$ up to Lagerl\"{o}f's calibration $\gamma=\gamma_c$. Notice how, at such critical value, the fixed point $L^*$ collides with the switching manifold $M_1$ (green curve, see Eq.~\eqref{eq:switching_manifold_1}) and a line of stable fixed points $\mathcal{X}^*$ (light blue line, see Eq.~\eqref{eq:line_fixed_points_asymp} is generated. This critical value identifies a border collision bifurcation for the fixed points $X^*$ and $\mathcal X^*$, which do not exist in $\mathbb{R}^2_{>0}$ for any $\gamma>\gamma_c$.
    Panels (b1)-(b3): state space $\mathbb{R}^2_{>0}$ of the asymptotic reduction at different slices of the bifurcation diagram: (b1) $\gamma=0.21<\gamma_c$; (b2) $\gamma=\gamma_c$ and; (b3) $\gamma=0.24>\gamma_c$.
    In each of those panels, the switching manifolds and their intersection are the same reported in Figure~\ref{fig:asym_phase}. 
    In addition we plot the fixed points $X^*$ (blue triangle) and $\mathcal X^*$ (light blue line) and the trajectory in state space (red curve) of an economy starting at initial condition $(G_0,L_0)=(1.0752,9.9109)\in\mathcal E$; notice that the $L_t$ component of each trajectory is shown in Figure~\ref{fig:different_steady_states}.
    }
    \label{fig:border_collision}
\end{figure}
%
This behavior can be entirely explained from a geometric argument pertaining to the persistence of these steady-states with changing values of $\gamma$.
%
\begin{proposition}\label{prop:border_collison}
    Let $X_t=(G_t,L_t)\in\mathbb{R}^2_{>0}$, $f_{\infty}$ given by \eqref{eq:asymptotic_reduction}-\eqref{eq:lagerlof_education}, $G^*$ given by \eqref{eq:fixed_points_asymp}, $\tau$, $\rho$ and $\theta$ given by \eqref{eq:lagerlof_parameters} and $\gamma=\gamma_c$, then
    %
    \begin{equation}\label{eq:line_fixed_points_asymp}
        \mathcal X^*=\bigg\{X_t\in\mathbb{R}^2_{>0}\;:\; G_t=G^*,\, L_t \geq \frac{\hat L}{\theta} \bigg\}\,,
    \end{equation}
    %
    is a set of stable fixed points of $f_{\infty}$ in $\mathbb{R}^2_{>0}$.
\end{proposition}
%
\begin{proof}
    This proof follows the same argument as the one outlined in Proposition~\ref{prop:asymp_fixed_points}. The only nuance here is to realize that, for $L_t>\frac{\hat L}{\theta}$ and $\gamma=\gamma_c$, imposing the fixed-point condition $L_{t+1}=L_t$ yields the identity $L^*=L^*$, which thus gives the set \eqref{eq:line_fixed_points_asymp}.
    \qed
\end{proof}
%
In fact, a conclusion of Propositions \ref{prop:asymp_fixed_points} and \ref{prop:border_collison} is that, at Lagerl\"{o}f's prescribed choice $\gamma=\gamma_c$, there is a qualitative change in the dynamics of economies in the modern-growth regime.

This is better illustrated in Figure~\ref{fig:border_collision}.
When $\gamma<\gamma_c$, Eq.~\eqref{eq:fixed_points_asymp} shows that the population component $L^*$ of the fixed point $X^*$ depends linearly on $\gamma$ (panel (a), blue line) and, crucially, it is smaller than the critical value $\hat L/\theta$ (panel (a), dashed green line) that identifies the switching manifold $M_1$ given in \eqref{eq:switching_manifold_1}.
If we let the weight parameter increase and approach Lagerl\"{o}f's value from below, $\gamma \uparrow \gamma_c$, we notice how the steady-state $X^*\in\mathbb R^2_{>0}$ (panel (b1), blue triangle) approaches the switching manifold $M_1$ (panel (b1), green line) from the region $\mathcal E$.
When the weight parameter matches exactly the one proposed by Lagerl\"{o}f, $\gamma = \gamma_c$, such steady-state collides with the switching manifold $M_1$ (panel (b2)) and disappears from the reduced state space.
In its place, a line of stable fixed points \eqref{eq:line_fixed_points_asymp} is created (panels (a)-(b3), light blue line).
Each of these fixed points correspond to possible steady-states of an economy in the $\mathcal D$ region of the asymptotic reduction.
Some of these steady-states are shown as timeseries in Figure~\ref{fig:different_steady_states}.
Finally, for any value of the weight parameter strictly greater than Lagerl\"{o}f choice, $\gamma > \gamma_c$, such steady-state no longer exists in the reduced state space $\mathbb R^2_{>0}$ (panel (b3)).

We conclude that, the asymptotic reduction $f_{\infty}$ is at a border collision bifurcation (recall Definition \textbf{INSERT REF.}) for the weight parameter's value $\gamma=\gamma_c$, and for the remaining parameter's values \eqref{eq:lagerlof_parameters}, specified by Lagerl\"{o}f.
In particular, Figure~\ref{fig:border_collision} is a bifurcation diagram of the sustained growth steady-state in the asymptotic reduction of the GWL model.

\section{An alternative formulation of the GWL model}\label{sec:smoothing}

\section{Conclusions}\label{sec:conclusions}

\printcredits

\section*{Declaration of competing interest}
The authors declare that they have no known competing financial interests or personal relationships that could have appeared to influence the work reported in this paper.

\section*{Data availability}
The code used to generate the figures is available at the repository \url{https://github.com/samuelbolduc44/UGT---codes/tree/main}.

%\bibliographystyle{cas-common}
\bibliographystyle{econ-aea}
\bibliography{references}

%%%%%%%%%%%%%%%%%%%%%%%%%%%%%%%%%%%%
%%%%%%   --- APPENDICES ---   %%%%%%
%%%%%%%%%%%%%%%%%%%%%%%%%%%%%%%%%%%%

\appendix
\section{Proof of Proposition 1}\label{appenix_A}
\samb{I am worried that the global existence proof is tricky and needs stronger conditions. }
\section{Proof of Lemma 2}\label{appenix_B}
Let us set $t=1$ in \eqref{eq:population_collapse_1} 
%
\begin{equation*}
        0\leq A_1<L_1\Big(\tilde{c}\,h_1^{-\alpha}\Big)^{\frac{1}{{1-\alpha}}}\,,
\end{equation*}
%
where $h_1=h(G_1)$ as per \eqref{eq:lagerlof_capital}.
We can then use the above inequality as a starting point to compute $U_2$ explicitly.
To do so we need to substitute the relevant identities for the state variables $(G_1,A_1,L_1)$ and the observables $e_1,h_1,z_1$ until we are left with a set of inequalities that only involves initial conditions of the state variables $(G_0,A_0,L_0)$.
Notice that this is done for the sake of simplicity and readability of the proof; the same result can be achieved by setting $t\to t-2$ in \eqref{eq:population_collapse_1} and solving it backwards to inequalities that only involves the state variables $(G_t,A_t,L_t)$.

We start by substituting the evolution law \eqref{eq:galor_family} for $A_1$ in the above inequality, as prescribed in \cite{Galor2000}.
After a trivial rearrangement, we get
%
\begin{equation}\label{eq:sufficient_condition}
        0\leq A_0<\frac{L_1\,\Big(\tilde{c}\,h_1^{-\alpha}\Big)^{\frac{1}{{1-\alpha}}}}{1 + G_1}\,,
\end{equation}
%
which is of the same form of \eqref{eq:population_collapse_1} but with high-order terms in the left-hand side term.

To prove that $U_2$ is disjoint from $U_1$ we need to set $L_1\neq0$. 
Therefore \eqref{eq:sufficient_condition} needs to be satisfied only for two of the three \emph{branches} of \eqref{eq:lagerlof_optimum}, specifically $L_1=n(G_0,\,A_0,\,L_0)L_0\neq0$ implies
%
\begin{equation}\label{eq:restricted_optimum}
    n(G_0,\,A_0,\,L_0) = 
    \begin{cases}
        \frac{\gamma}{\tau + e_{1}}, & \frac{\tilde c}{1-\gamma} \leq z_0, \\[8pt]
        \frac{1 - \frac{\tilde c}{z_0}}{\tau + e_{1}}, & \tilde c \leq z_0 < \frac{\tilde c}{1-\gamma}. \\[8pt]
    \end{cases}
\end{equation}
%
Combining \eqref{eq:restricted_optimum} with \eqref{eq:sufficient_condition} we conclude that the set $U_2$ is made of initial conditions $(G_0,A_0,L_0)$ that must satisfy both
%
\begin{equation}\label{eq:upper}
    A_0 < \frac{\gamma\,L_0\,\Big(\tilde{c}\,h_1^{-\alpha}\Big)^{\frac{1}{{1-\alpha}}}}{\big(\tau + e_1\big)\big(1 + G_1\big)}\,,
\end{equation}
%
and
\begin{equation}\label{eq:lower}
    A_0 < \frac{\Big(1-\frac{\tilde{c}}{z_0}\Big)\,L_0\,\Big(\tilde{c}\,h_1^{-\alpha}\Big)^{\frac{1}{{1-\alpha}}}}{\big(\tau + e_1\big)\big(1 + (e_0 + \tau\rho)\min(\theta L_0,\,\hat{L})\big)}\,,
\end{equation}
%
where $G_1$ is given by the evolution law \eqref{eq:lagerlof_choices}
\begin{equation}\label{eq:accessory_G1}
    G_1 = (e_0 + \tau\rho)\min(\theta L_0,\hat{L})\,,
\end{equation}
%
as prescribed in \cite{Lagerlof2006}, and the observables $e_{0,1}=e(G_{0,1})$, $h_1=h(G_1)$ and $z_0=z(G_0,A_0,L_0)$ are given by \eqref{eq:lagerlof_education}, \eqref{eq:lagerlof_capital} and \eqref{eq:lagerlof_income} respectively
%
\begin{subequations}\label{eq:accessories}
\begin{align}
    e_0 &= \max\!\Big(0,\,-\rho\tau+\sqrt{\tau(1-\rho)G_0}\Big), \label{eq:accessory_e0}\\
    e_1 &= \max\!\Big(0,\,-\rho\tau+\sqrt{\tau(1-\rho)G_1}\Big), \label{eq:accessory_e1}\\
    h_1 &= \frac{e_1 + \rho\tau}{e_1 + \rho\tau + G_1}, \label{eq:accessory_h1}\\
    z_0 &= \bigg(\frac{e_0 + \rho\tau}{e_0 + \rho\tau + G_0}\bigg)^{\alpha}\bigg(\frac{A_0}{L_0}\bigg)^{1-\alpha}. \label{eq:accessory_z0}
\end{align}
\end{subequations}
%

Notice that the switching terms $\min(\theta L_0,\hat{L})$ and $\max(0,\cdot)$ partition the subspace $(G_0,\,L_0)$ into distinct regions for which \eqref{eq:upper} and \eqref{eq:lower} hold:
%
\begin{enumerate}
    \item from \eqref{eq:accessory_G1} the $L_0-$coordinate is partitioned in
    \begin{subequations}\label{eq:partition_L0}
    \begin{align}
        0 &\leq L_0 < \frac{\hat{L}}{\theta}\quad\implies\quad G_1=(e_0+\tau\rho)\theta L_0\,, \label{eq:partition_L0_1} \\ 
        \frac{\hat{L}}{\theta} &\leq L_0 \quad\implies\quad G_1=(e_0+\tau\rho)\hat{L}\,; \label{eq:partition_L0_2}
    \end{align}
    \end{subequations}
    %
    \item from \eqref{eq:accessory_e0} the $G_0-$coordinate is partitioned in
    \begin{subequations}\label{eq:partition_G0}
    \begin{align}
        0 &\leq G_0 < \frac{\rho^2\tau}{1-\rho}\quad\implies\quad e_0=0\,, \label{eq:partition_G0_1} \\ 
        \frac{\rho^2\tau}{1-\rho} &\leq G_0 \quad\implies\quad e_0=-\rho\tau+\sqrt{\tau(1-\rho)G_0}\,; \label{eq:partition_G0_2}
    \end{align}
    \end{subequations}
    %
    \item from \eqref{eq:accessory_e1} we observe a nonlinear coupling of the $L_0-$ and $G_0-$coordinate
    %
    \begin{equation*}
        \bigg(\max\Big(0,-\rho\tau + \sqrt{\tau(1-\rho)G_0}\Big) + \tau\rho\bigg)\min(\theta L_0,\hat{L}) = \frac{\rho^2\tau}{1-\rho}\,.
    \end{equation*}
    %
    From \eqref{eq:partition_L0} and \eqref{eq:partition_G0} we have
    %
    \begin{subequations}\label{eq:partition_L0_G0}
    \begin{align}
        0 &\leq L_0 < \frac{\hat{L}}{\theta}\quad\text{and}\quad0\leq G_0 < \frac{\rho^2\tau}{1-\rho}\quad \implies\quad e_1=\tau\,\max\!\Big(0,\,-\rho+\sqrt{(1-\rho)\rho\,\theta L_0}\Big)\,, \label{eq:partition_L0_G0_1} \\
        %
        0 &\leq L_0 < \frac{\hat{L}}{\theta}\quad\text{and}\quad\frac{\rho^2\tau}{1-\rho}\leq G_0\; \implies\; e_1=\max\!\bigg(0,\,-\rho\tau+\sqrt{\tau(1-\rho)\sqrt{\tau(1-\rho)G_0}\,\theta L_0}\bigg)\,, \label{eq:partition_L0_G0_2} \\
        %
        \frac{\hat{L}}{\theta} &\leq L_0\quad\text{and}\quad0\leq G_0 < \frac{\rho^2\tau}{1-\rho}\quad \implies\quad e_1=\tau\,\max\!\Big(0,\,-\rho+\sqrt{(1-\rho)\rho\,\hat{L}}\Big)\,, \label{eq:partition_L0_G0_3} \\
        %
        \frac{\hat{L}}{\theta} &\leq L_0\quad\text{and}\quad\frac{\rho^2\tau}{1-\rho}\leq G_0\quad \implies\quad e_1=\max\!\bigg(0,\,-\rho\tau+\sqrt{\tau(1-\rho)\sqrt{\tau(1-\rho)G_0}\,\hat{L}}\bigg) \label{eq:partition_L0_G0_4}\,.
    \end{align}
    \end{subequations}
    %
\end{enumerate}
%
Given Lagerl\"{o}f's choice of parameters $\tau=0.15$, $\rho=0.879$, $\alpha=0.6$, $\theta=1$ and $\hat{L}=11.42$, we combine conditions \eqref{eq:partition_L0_G0} with \eqref{eq:partition_L0} and \eqref{eq:partition_G0} to get the partition of the $(G_0,L_0)-$plane:
%
\begin{enumerate}
    \item from \eqref{eq:partition_L0_G0_1} the $(G_0,L_0)-$plane is partitioned in
    %
    \begin{subequations}\label{eq:sub_partition_1}
    \begin{align}
        &0 \leq L_0 < \frac{\rho}{(1-\rho)\theta}\quad\text{and}\quad0\leq G_0 < \frac{\rho^2\tau}{1-\rho}\quad\implies\quad 
        \begin{aligned}[t]
            e_0 &= 0, \\
            e_1 &= 0,\\
            G_1 &= \tau\rho\,\theta L_0,
        \end{aligned}
        \label{eq:sub_partition_1_1}\\
        %
        &\frac{\rho}{(1-\rho)\theta} \leq L_0 < \frac{\hat{L}}{\theta}\quad\text{and}\quad0\leq G_0 < \frac{\rho^2\tau}{1-\rho}\quad\implies\quad
        \begin{aligned}[t]
            e_0 &= 0, \\
            e_1 &= \tau\Big(\sqrt{(1-\rho)\rho\,\theta L_0}-\rho\Big),\\
            G_1 &= \tau\rho\,\theta L_0\,;
        \end{aligned}
        \label{eq:sub_partition_1_2}
    \end{align}
    \end{subequations}
    %
    \item from \eqref{eq:partition_L0_G0_2} the $(G_0,L_0)-$plane is partitioned in
    %
    \begin{subequations}\label{eq:sub_partition_2}
    \begin{align}
        &0 \leq L_0 < \frac{\hat{L}}{\theta}\quad\text{and}\quad\frac{\rho^2\tau}{1-\rho}\leq G_0 < \frac{\rho^4\tau}{(1-\rho)^3\,\theta^2}\,L_0^{-2}\quad\implies\quad 
        \begin{aligned}[t]
            e_0 &= -\rho\tau + \sqrt{\tau(1-\rho)G_0}, \\
            e_1 &= 0,\\
            G_1 &= \theta L_0\sqrt{\tau(1-\rho)G_0},
        \end{aligned}
        \label{eq:sub_partition_2_1}\\
        %
        &0 \leq L_0 < \frac{\rho}{(1-\rho)\tau}\quad\text{and}\quad\frac{\rho^4\tau}{(1-\rho)^3\,\theta^2}\,L_0^{-2}\leq G_0\;\implies\; 
        \begin{aligned}[t]
            e_0 &= -\rho\tau + \sqrt{\tau(1-\rho)G_0}, \\
            e_1 &= -\rho\tau+\sqrt{\tau(1-\rho)\sqrt{\tau(1-\rho)G_0}\,\theta L_0},\\
            G_1 &= \theta L_0\sqrt{\tau(1-\rho)G_0}\,,
        \end{aligned}
        \label{eq:sub_partition_2_2}\\
        %
        &\frac{\rho}{(1-\rho)\tau}\leq L_0 < \frac{\hat{L}}{\theta}\quad\text{and}\quad\frac{\rho^2\tau}{1-\rho}\leq G_0\;\implies\; 
        \begin{aligned}[t]
            e_0 &= -\rho\tau + \sqrt{\tau(1-\rho)G_0}, \\
            e_1 &= -\rho\tau+\sqrt{\tau(1-\rho)\sqrt{\tau(1-\rho)G_0}\,\theta L_0},\\
            G_1 &= \theta L_0\sqrt{\tau(1-\rho)G_0}\,,
        \end{aligned}
        \label{eq:sub_partition_2_3}
        %
    \end{align}
    \end{subequations}
    where we remark that \eqref{eq:sub_partition_2_2} and \eqref{eq:sub_partition_2_3} are indistinguishable for their values of $e_0$, $e_1$ and $G_1$ and therefore we can write it as
    %
    \begin{equation}\label{eq:sub_partition_2_2_and_3}
        0 \leq L_0 < \frac{\hat{L}}{\theta}\quad\text{and}\quad\max\!\left(\frac{\rho^2\tau}{1-\rho},\,\frac{\rho^4\tau}{(1-\rho)^3\,\theta^2}\,L_0^{-2}\right)\leq G_0\;\implies\; \eqref{eq:sub_partition_2_2}\equiv\eqref{eq:sub_partition_2_3}
    \end{equation}
    %
    \item from \eqref{eq:partition_L0_G0_3} no further partition is enforced since $e_1$ does not depend on $G_0$ nor $L_0$
    %
    \begin{equation}\label{eq:sub_partition_3}
        \frac{\hat{L}}{\theta} \leq L_0\quad\text{and}\quad0\leq G_0 < \frac{\rho^2\tau}{1-\rho}\quad \implies\quad
        \begin{aligned}[t]
            e_0 &= 0, \\
            e_1 &= \tau\Big(\sqrt{(1-\rho)\rho\,\hat{L}}-\rho\Big),\\
            G_1 &= \tau\rho\hat{L}\,;
        \end{aligned}
    \end{equation}
    %
    \item from \eqref{eq:partition_L0_G0_4} no further partition is enforced since $\frac{\rho^4\tau}{(1-\rho)^3\hat{L}^2} < \frac{\rho^2\tau}{1-\rho}$
    %
    \begin{equation}\label{eq:sub_partition_4}
        \frac{\hat{L}}{\theta}\leq L_0\quad\text{and}\quad\frac{\rho^2\tau}{1-\rho}\leq G_0\quad\implies\quad
        \begin{aligned}[t]
            e_0 &= -\rho\tau+\sqrt{\tau(1-\rho)G_0},\\
            e_1 &= -\rho\tau+\sqrt{\tau(1-\rho)\sqrt{\tau(1-\rho)G_0}\hat{L}},\\
            G_1 &= \sqrt{\tau(1-\rho)G_0}\,\hat{L},
        \end{aligned}
    \end{equation}
\end{enumerate}
%

It follows that the subset $U_2$ is obtained by evaluating \eqref{eq:upper} and \eqref{eq:lower} for the six partitions of the $(G_0,L_0)-$plane given in \eqref{eq:sub_partition_1}-\eqref{eq:sub_partition_4}.
Let us denote $\Omega\coloneqq\mathbb{R}^2_{\geq0}$ to be the subset of the $(G_0,L_0)-$plane in which $f$ is defined.
Based on the above, we introduce the following partition of $\Omega$
%
\begin{subequations}\label{eq:partition_omega}
    \begin{align}
        \Omega_1&\coloneqq\bigg\{(G_0,L_0)\in\Omega\;:\; 0 \leq G_0 < \frac{\rho^2\tau}{1-\rho}\,, \;0 \leq L_0 < \frac{\rho}{(1-\rho)\theta}\bigg\}\,,\label{eq:partition_omega_1}\\
        %
        \Omega_2&\coloneqq\bigg\{0 \leq G_0 < \frac{\rho^2\tau}{1-\rho}\,, \;\frac{\rho}{(1-\rho)\theta} \leq L_0 < \frac{\hat{L}}{\theta}\bigg\}\,,\label{eq:partition_omega_2}\\
        %
        \Omega_3&\coloneqq\bigg\{0 \leq G_0 < \frac{\rho^2\tau}{1-\rho} \,, \; \frac{\hat{L}}{\theta}\leq L_0\bigg\}\,,\label{eq:partition_omega_3}\\
        %
        \Omega_4&\coloneqq\bigg\{\frac{\rho^2\tau}{1-\rho} \leq G_0 < \frac{\rho^4\tau}{(1-\rho)^3\,\theta^2}\,L_0^{-2} \,, \;0 \leq L_0 < \frac{\rho}{(1-\rho)\theta}\bigg\}\,,\label{eq:partition_omega_4}\\
        %
        \Omega_5&\coloneqq\bigg\{\frac{\rho^4\tau}{(1-\rho)^3\,\theta^2}\,L_0^{-2} \leq G_0 \,, \; 0 \leq L_0 < \frac{\rho}{(1-\rho)\theta}\bigg\} \bigcup \bigg\{\frac{\rho^2\tau}{1-\rho} \leq G_0 \,, \; \frac{\rho}{(1-\rho)\theta} \leq L_0 < \frac{\hat{L}}{\theta}\bigg\}\,,\label{eq:partition_omega_5}\\
        %
        \Omega_6&\coloneqq\bigg\{\frac{\rho^2\tau}{1-\rho} \leq G_0 \,, \; \frac{\hat{L}}{\theta}\leq L_0\bigg\}\,.\label{eq:partition_omega_6}
    \end{align}
\end{subequations}
%
By denoting $U_2^{(j)}$ to be the subset of $U_2$ obtained by those $(G_0,A_0,L_0)$ such that \eqref{eq:upper} and \eqref{eq:lower} hold for all $(G_0,L_0)\in\Omega_j$, we construct $U_2=\bigcup_{j=1}^6 U_2^{(j)}$.
This is straightforward to achieve since each $\Omega_j$ maps one-to-one to a triple $(e_0,e_1,G_1)$ as follows
%
\begin{subequations}\label{eq:partition_mapping}
    \begin{align}
        \eqref{eq:partition_omega_1}& \quad\longrightarrow\quad \eqref{eq:sub_partition_1_1}\,,\label{eq:partition_mapping_1}\\
        %
        \eqref{eq:partition_omega_2}& \quad\longrightarrow\quad \eqref{eq:sub_partition_1_2}\,,\label{eq:partition_mapping_2}\\
        %
        \eqref{eq:partition_omega_3}& \quad\longrightarrow\quad \eqref{eq:sub_partition_3}\,,\label{eq:partition_mapping_3}\\
        %
        \eqref{eq:partition_omega_4}& \quad\longrightarrow\quad \eqref{eq:sub_partition_2_1}\,,\label{eq:partition_mapping_4}\\
        %
        \eqref{eq:partition_omega_5}& \quad\longrightarrow\quad \eqref{eq:sub_partition_2_2_and_3}\,,\label{eq:partition_mapping_5}\\
        %
        \eqref{eq:partition_omega_6}& \quad\longrightarrow\quad \eqref{eq:sub_partition_4}\,.\label{eq:partition_mapping_6}
    \end{align}
\end{subequations}
%

What is left is the evaluation of \eqref{eq:upper} and \eqref{eq:lower} using each set of equations in \eqref{eq:partition_mapping} to build each $U_2$ piece-by-piece (i.e. assembling each $U_2^{(j)}$ individually).
The computations are as follows:
%
\begin{enumerate}
    \item from \eqref{eq:partition_mapping_1} we get, $\forall (G_0,L_0)\in\Omega_1$,
    \begin{equation}\label{eq:subset_M2_1}
        U_2^{(1)} = \left\{
        \begin{aligned}
            A_0 &< \frac{\gamma\,L_0\,\Big[\tilde{c}\,(1+\theta L_0)^{\alpha}\Big]^{\frac{1}{1-\alpha}}}{\tau\big(1+\tau\rho\,\theta L_0\big)}\,,\\[10pt]
            %
            A_0 &< \frac{L_0\,\Big[\tilde{c}\,(1+\theta L_0)^{\alpha}\Big]^{\frac{1}{1-\alpha}}\left(1-\tilde{c}\left(\frac{\rho\tau+G_0}{\rho\tau}\right)^{\alpha}\left(\frac{L_0}{A_0}\right)^{1-\alpha}\right)}{\tau\big(1+\tau\rho\,\theta L_0\big)}
        \end{aligned}
        \right\}\,;
    \end{equation}
    %
    \item from \eqref{eq:partition_mapping_2} we get, $\forall (G_0,L_0)\in\Omega_2$,
    \begin{equation}\label{eq:subset_M2_2}
        U_2^{(2)} = \left\{
        \begin{aligned}
            A_0 &< \frac{\gamma\,L_0\,\left[\tilde{c}\,\left(\dfrac{\sqrt{(1-\rho)\rho\,\theta L_0} - \rho(1-\tau) + \tau\rho\,\theta L_0}{\sqrt{(1-\rho)\rho\,\theta L_0} - \rho(1-\tau)}\right)^{\alpha}\right]^{\frac{1}{1-\alpha}}}{\Big(\tau-\rho+\sqrt{(1-\rho)\rho\,\theta L_0}\Big)\big(1+\tau\rho\,\theta L_0\big)}\,,\\[10pt]
            %
            A_0 &< \frac{L_0\,\left[\tilde{c}\,\left(\dfrac{\sqrt{(1-\rho)\rho\,\theta L_0} - \rho(1-\tau) + \tau\rho\,\theta L_0}{\sqrt{(1-\rho)\rho\,\theta L_0} - \rho(1-\tau)}\right)^{\alpha}\right]^{\frac{1}{1-\alpha}}\left(1-\tilde{c}\left(\frac{\rho\tau+G_0}{\rho\tau}\right)^{\alpha}\left(\frac{L_0}{A_0}\right)^{1-\alpha}\right)}{\Big(\tau-\rho+\sqrt{(1-\rho)\rho\,\theta L_0}\Big)\big(1+\tau\rho\,\theta L_0\big)}
        \end{aligned}
            \right\}\,;
    \end{equation}
    %
    \item from \eqref{eq:partition_mapping_3} we get, $\forall (G_0,L_0)\in\Omega_3$,
    \begin{equation}\label{eq:subset_M2_3}
        U_2^{(3)} = \left\{ 
        \begin{aligned}
            A_0 &< \frac{\gamma\,L_0\,\left[\tilde{c}\,\left(1+\sqrt{\dfrac{\rho\hat{L}}{1-\rho}}\right)^{\alpha}\right]^{\frac{1}{1-\alpha}}}{\tau\Big(1-\rho+\sqrt{(1-\rho)\rho\,\hat{L}}\Big)\big(1+\tau\rho\,\hat{L}\big)}\,,\\[10pt]
            %
            A_0 &< \frac{L_0\,\left[\tilde{c}\,\left(1+\sqrt{\dfrac{\rho\hat{L}}{1-\rho}}\right)^{\alpha}\right]^{\frac{1}{1-\alpha}}\left(1-\tilde{c}\left(\frac{\rho\tau+G_0}{\rho\tau}\right)^{\alpha}\left(\frac{L_0}{A_0}\right)^{1-\alpha}\right)}{\tau\Big(1-\rho+\sqrt{(1-\rho)\rho\,\hat{L}}\Big)\big(1+\tau\rho\,\hat{L}\big)}
        \end{aligned}
        \right\}\,;
    \end{equation}
    %
    \item from \eqref{eq:partition_mapping_4} we get, $\forall (G_0,L_0)\in\Omega_4$,
    \begin{equation}\label{eq:subset_M2_4}
        U_2^{(4)} = \left\{ 
        \begin{aligned}
            A_0 &< \frac{\gamma\,L_0\,\left[\tilde{c}\,\left(1+\dfrac{\theta L_0\sqrt{\tau(1-\rho)G_0}}{\rho\tau}\right)^{\alpha}\right]^{\frac{1}{1-\alpha}}}{\tau\Big(1+\theta L_0\sqrt{\tau(1-\rho)G_0}\Big)}\,,\\[10pt]
            %
            A_0 &< \frac{L_0\,\left[\tilde{c}\,\left(1+\dfrac{\theta L_0\sqrt{\tau(1-\rho)G_0}}{\rho\tau}\right)^{\alpha}\right]^{\frac{1}{1-\alpha}}\left(1-\tilde{c}\left(\frac{\sqrt{\tau(1-\rho)G_0}+G_0}{\sqrt{\tau(1-\rho)G_0}}\right)^{\alpha}\left(\frac{L_0}{A_0}\right)^{1-\alpha}\right)}{\tau\Big(1+\theta L_0\sqrt{\tau(1-\rho)G_0}\Big)}
        \end{aligned}
        \right\}\,;
    \end{equation}
    %
    \item from \eqref{eq:partition_mapping_5} we get, $\forall (G_0,L_0)\in\Omega_5$,
    \begin{equation}\label{eq:subset_M2_5}
        U_2^{(5)} = \left\{ 
        \begin{aligned}
            A_0 &< \frac{\gamma\,L_0\,\left[\tilde{c}\,\left(1+\dfrac{\theta L_0\sqrt{\tau(1-\rho)G_0}}{\sqrt{\tau(1-\rho)\,\theta L_0\sqrt{\tau(1-\rho)G_0}}}\right)^{\alpha}\right]^{\frac{1}{1-\alpha}}}{\left(\tau(1-\rho)+\sqrt{\tau(1-\rho)\,\theta L_0\sqrt{\tau(1-\rho)G_0}}\right)\Big(1+\theta L_0\sqrt{\tau(1-\rho)G_0}\Big)}\,,\\[10pt]
            %
            A_0 &< \frac{L_0\,\left[\tilde{c}\,\left(1+\dfrac{\theta L_0\sqrt{\tau(1-\rho)G_0}}{\sqrt{\tau(1-\rho)\,\theta L_0\sqrt{\tau(1-\rho)G_0}}}\right)^{\alpha}\right]^{\frac{1}{1-\alpha}}\left(1-\tilde{c}\left(\frac{\sqrt{\tau(1-\rho)G_0}+G_0}{\sqrt{\tau(1-\rho)G_0}}\right)^{\alpha}\left(\frac{L_0}{A_0}\right)^{1-\alpha}\right)}{\left(\tau(1-\rho)+\sqrt{\tau(1-\rho)\,\theta L_0\sqrt{\tau(1-\rho)G_0}}\right)\Big(1+\theta L_0\sqrt{\tau(1-\rho)G_0}\Big)}
        \end{aligned}
        \right\}\,;
    \end{equation}
    %
    \item from \eqref{eq:partition_mapping_6} we get, $\forall (G_0,L_0)\in\Omega_6$,
    \begin{equation}\label{eq:subset_M2_6}
        U_2^{(6)} = \left\{ 
        \begin{aligned}
            A_0 &< \frac{\gamma\,L_0\,\left[\tilde{c}\,\left(1+\dfrac{\hat{L}\sqrt{\tau(1-\rho)G_0}}{\sqrt{\tau(1-\rho)\,\hat{L}\sqrt{\tau(1-\rho)G_0}}}\right)^{\alpha}\right]^{\frac{1}{1-\alpha}}}{\left(\tau(1-\rho)+\sqrt{\tau(1-\rho)\,\hat{L}\sqrt{\tau(1-\rho)G_0}}\right)\Big(1+\hat{L}\sqrt{\tau(1-\rho)G_0}\Big)}\,,\\[10pt]
            %
            A_0 &< \frac{L_0\,\left[\tilde{c}\,\left(1+\dfrac{\hat{L}\sqrt{\tau(1-\rho)G_0}}{\sqrt{\tau(1-\rho)\,\hat{L}\sqrt{\tau(1-\rho)G_0}}}\right)^{\alpha}\right]^{\frac{1}{1-\alpha}}\left(1-\tilde{c}\left(\frac{\sqrt{\tau(1-\rho)G_0}+G_0}{\sqrt{\tau(1-\rho)G_0}}\right)^{\alpha}\left(\frac{L_0}{A_0}\right)^{1-\alpha}\right)}{\left(\tau(1-\rho)+\sqrt{\tau(1-\rho)\,\hat{L}\sqrt{\tau(1-\rho)G_0}}\right)\Big(1+\hat{L}\sqrt{\tau(1-\rho)G_0}\Big)}
        \end{aligned}
        \right\}\,.
    \end{equation}
\end{enumerate}

From \eqref{eq:subset_M2_1}-\eqref{eq:subset_M2_6}, the proof follows.
\qed

\section{Dynamics and transition graph in the modern-growth regime}\label{appenix_C}

We give explicit formulations of the dynamical system in the regions $\mathcal A$--$\mathcal F$ for the asymptotic reduction and we show how it induces a directed graph $\Gamma$ that explains the transitions among those regions.

\subsection{Explicit dynamics in $\mathcal A$--$\mathcal F$}\label{appenix_C1}

We derive the explicit form of $f_{\infty}$, given by \eqref{eq:asymptotic_reduction}-\eqref{eq:lagerlof_education}, in the six different regions $\mathcal A$--$\mathcal F$ given by the switching manifolds \eqref{eq:switching_manifolds}.
For convenience, we explicitly report the dynamics in the asymptotic reduction \eqref{eq:asymptotic_reduction}
%
\begin{subequations}\label{eq:system_c}
    \begin{align}
        G_{t+1} &= (e(G_t) + \rho\tau)\min(\theta L_t, \hat L)\,, \label{eq:growth_c} \\
        L_{t+1} &= \frac{\gamma}{\tau + e(G_{t+1})} L_t\,, \label{eq:population_c}
    \end{align}
\end{subequations}
%
as well as the functional form for the choice variable education \eqref{eq:lagerlof_education}
%
\begin{equation}\label{eq:education_c}
    e(G_t)=\max\!\big(0,-\rho\tau + \sqrt{\tau(1-\rho)G_t}\big)\,,
\end{equation}
%
and proceed with the derivation of those rules in each region.

\subsubsection{Region $\mathcal{A}$}\label{appenix_C_1}

In this region, $G_t<\frac{\rho^2\tau}{1-\rho}$ and $\frac{\hat{L}}{\theta}<L_t$. 
From \eqref{eq:education_c} this implies $e(G_t) = 0$, which, from \eqref{eq:growth_c}, gives
%
\begin{equation*}
    G_{t+1} = \rho\tau\hat L\,.
\end{equation*}
%
Plugging the above back in \eqref{eq:education_c}, and solving for $L_t$, gives the dynamics in region $\mathcal A$
%
\begin{equation}\label{eq:dynamics_in_A}
    f_{\infty}(X_t\in\mathcal A) = \begin{cases}
                            G_{t+1} = \rho\tau\hat L_t\,, \\[8pt]
                            L_{t+1} = \frac{\gamma}{\tau(1-\rho) + \sqrt{\hat{L}(1-\rho)\rho}} L_t\,.
                      \end{cases}
\end{equation}

\subsubsection{Region $\mathcal{B}$}\label{appenix_C_2}

In this region, $G_t<\frac{\rho^2\tau}{1-\rho}$ and $\frac{\rho}{\theta(1-\rho)}<L_t<\frac{\hat{L}}{\theta}$. 
From \eqref{eq:education_c} this implies $e(G_t) = 0$, which, from \eqref{eq:growth_c}, gives
%
\begin{equation*}
    G_{t+1} = \rho\tau\theta L_t\,.
\end{equation*}
%
Plugging the above back in \eqref{eq:education_c}, and solving for $L_t$, gives the dynamics in region $\mathcal B$
%
\begin{equation}\label{eq:dynamics_in_B}
    f_{\infty}(X_t\in\mathcal B) = \begin{cases}
                            G_{t+1} = \rho\tau\theta L_t\,, \\[8pt]
                            L_{t+1} = \frac{\gamma}{\tau(1-\rho) + \tau\sqrt{(1-\rho)\rho\theta L_t}} L_t\,.
                      \end{cases}
\end{equation}

\subsubsection{Region $\mathcal{C}$}\label{appenix_C_3}

In this region, $G_t<\frac{\rho^2\tau}{1-\rho}$ and $L_t<\frac{\rho}{\theta(1-\rho)}$. 
From \eqref{eq:education_c} this implies $e(G_t) = 0$, which, from \eqref{eq:growth_c}, gives
%
\begin{equation*}
    G_{t+1} = \rho\tau\theta L_t\,.
\end{equation*}
%
Plugging the above back in \eqref{eq:education_c}, and solving for $L_t$, gives the dynamics in region $\mathcal C$
%
\begin{equation}\label{eq:dynamics_in_C}
    f_{\infty}(X_t\in\mathcal C) = \begin{cases}
                            G_{t+1} = \rho\tau\theta L_t\,, \\[8pt]
                            L_{t+1} = \frac{\gamma}{\tau} L_t\,.
                      \end{cases}
\end{equation}

\subsubsection{Region $\mathcal{D}$}\label{appenix_C_4}

In this region, $\frac{\rho^2\tau}{1-\rho} < G_t$ and $\frac{\hat L}{\theta}<L_t$. 
From \eqref{eq:education_c} this implies $e(G_t) = - \rho\tau + \sqrt{\tau(1-\rho)G_t}$, which, from \eqref{eq:growth_c}, gives
%
\begin{equation*}
    G_{t+1} = \hat L\sqrt{\tau(1-\rho) G_t}\,.
\end{equation*}
%
Plugging the above back in \eqref{eq:education_c}, and solving for $L_t$, gives the dynamics in region $\mathcal D$
%
\begin{equation}\label{eq:dynamics_in_D}
    f_{\infty}(X_t\in\mathcal D) = \begin{cases}
                            G_{t+1} = \hat L\sqrt{\tau(1-\rho)G_t} \,, \\[8pt]
                            L_{t+1} = \frac{\gamma}{\tau(1-\rho) + \hat L\sqrt{\tau(1-\rho)G_t}} L_t\,.
                      \end{cases}
\end{equation}

\subsubsection{Region $\mathcal{E}$}\label{appenix_C_5}

In this region, $\frac{\rho^2\tau}{1-\rho} < G_t$ and $\frac{\rho}{\theta(1-\rho)}<L_t<\frac{\hat L}{\theta}$. 
From \eqref{eq:education_c} this implies $e(G_t) = - \rho\tau + \sqrt{\tau(1-\rho)G_t}$, which, from \eqref{eq:growth_c}, gives
%
\begin{equation*}
    G_{t+1} = \theta L_t\sqrt{\tau(1-\rho) G_t}\,.
\end{equation*}
%
Plugging the above back in \eqref{eq:education_c}, and solving for $L_t$, gives the dynamics in region $\mathcal E$
%
\begin{equation}\label{eq:dynamics_in_E}
    f_{\infty}(X_t\in\mathcal E) = \begin{cases}
                            G_{t+1} = \theta L_t\sqrt{\tau(1-\rho)G_t} \,, \\[8pt]
                            L_{t+1} = \frac{\gamma}{\tau(1-\rho) + \theta L_t\sqrt{\tau(1-\rho)G_t}} L_t\,.
                      \end{cases}
\end{equation}

\subsubsection{Region $\mathcal{F}$}\label{appenix_C_6}

In this region, $\frac{\rho^2\tau}{1-\rho} < G_t$ and $L_t<\frac{\rho}{\theta(1-\rho)}$. 
From \eqref{eq:education_c} this implies $e(G_t) = - \rho\tau + \sqrt{\tau(1-\rho)G_t}$, which, from \eqref{eq:growth_c}, gives
%
\begin{equation*}
    G_{t+1} = \theta L_t\sqrt{\tau(1-\rho) G_t}\,.
\end{equation*}
%
Plugging the above back in \eqref{eq:education_c}, and solving for $L_t$, gives the dynamics in region $\mathcal F$
%
\begin{equation}\label{eq:dynamics_in_F}
    f_{\infty}(X_t\in\mathcal F) = \begin{cases}
                            G_{t+1} = \theta L_t\sqrt{\tau(1-\rho)G_t} \,, \\[8pt]
                            L_{t+1} = \frac{\gamma}{\tau} L_t\,.
                      \end{cases}
\end{equation}

\subsection{Construction of $\Gamma$}\label{appenix_C2}

\end{document}